\chapter{Chapter07 — Hastad Three Sat}
%%% generated by the blueprint pipeline from FABL.Chapter07.HastadThreeSat
\section{Overview}

Declarations of the Lean module \texttt{FABL.Chapter07.HastadThreeSat}.

\section{Declarations}

\begin{definition}[Predicate indicator function]
\lean{FABL.SignPredicate.indicator}
\statementsource{Analysis-of-Boolean-Functions-by-Ryan-ODonnell}{pdf-ffbe9c7b3306-p013-b006}
\label{FABL.SignPredicate.indicator}
\leanok
\uses{FABL.SignPredicate}
Let $P$ be a predicate on sign vectors of arity $k$. Define its indicator on the
$k$-dimensional sign cube to equal $1$ at inputs accepted by $P$ and $0$ at inputs
rejected by $P$.
\end{definition}

\begin{definition}[Multilinear extension of a predicate]
\lean{FABL.SignPredicate.multilinearExtension}
\statementsource{Analysis-of-Boolean-Functions-by-Ryan-ODonnell}{
  pdf-ffbe9c7b3306-p022-b002,
  pdf-ffbe9c7b3306-p035-b004
}
\label{FABL.SignPredicate.multilinearExtension}
\leanok
\uses{FABL.SignPredicate, FABL.SignPredicate.indicator, FABL.fourierCoeff}
Let $P$ be a predicate of arity $k$, and assign a real number $\mu_i$ to each input
position $i$. Define its multilinear extension to be the sum, over all subsets $S$ of
input positions, of the Fourier coefficient of the indicator of $P$ at $S$ multiplied by
$\prod_{i\in S}\mu_i$.
\end{definition}

\begin{theorem}[Expectation of a product under independent sampling]
\lean{FABL.pmfExpectation_independentProduct_subsetProd}
\label{FABL.pmfExpectation_independentProduct_subsetProd}
\leanok
\uses{FABL.independentProductPMF, FABL.pmfExpectation, FABL.pmfExpectation_const_one, FABL.pmfExpectation_independentProductPMF_prod}
\difficulty{4}
Let $I$ be a finite index set, let $A_i$ be a finite set with probability mass function
$p_i$ for each $i\in I$, and let $q_i:A_i\to\mathbb{R}$. For any subset $S\subseteq I$,
the expectation under the independent product distribution of $\prod_{i\in S}q_i(X_i)$
equals $\prod_{i\in S}\mathbb{E}_{X_i\sim p_i}[q_i(X_i)]$.
\end{theorem}

\begin{theorem}[Sign monomial expectation under independent sampling]
\lean{FABL.pmfExpectation_independentSigns_monomial}
\label{FABL.pmfExpectation_independentSigns_monomial}
\leanok
\uses{FABL.Sign, FABL.independentProductPMF, FABL.monomial, FABL.pmfExpectation, FABL.pmfExpectation_const_one, FABL.pmfExpectation_independentProductPMF_prod, FABL.pmfExpectation_signPMFOfMean, FABL.signPMFOfMean, FABL.signValue}
\difficulty{4}
Let $n\geq 0$, let $\mu_i\in[-1,1]$ for each $i\in\{0,\ldots,n-1\}$, and let
$S\subseteq\{0,\ldots,n-1\}$. Choose independent signs $X_i\in\{-1,1\}$ with mean
$\mu_i$. Then the expectation of the monomial $\prod_{i\in S}X_i$ is $\prod_{i\in
S}\mu_i$.
\end{theorem}

\begin{theorem}[Randomized value and multilinear extension]
\lean{FABL.SignPredicate.randomizedValue_eq_multilinearExtension}
\label{FABL.SignPredicate.randomizedValue_eq_multilinearExtension}
\leanok
\uses{FABL.SignPredicate, FABL.SignPredicate.indicator, FABL.SignPredicate.multilinearExtension, FABL.SignPredicate.randomizedValue, FABL.fourierCoeff, FABL.fourier_expansion, FABL.independentProductPMF, FABL.monomial, FABL.pmfExpectation, FABL.pmfExpectation_const_mul, FABL.pmfExpectation_independentSigns_monomial, FABL.pmfExpectation_sum, FABL.signPMFOfMean}
\difficulty{4}
Let $P$ be a Boolean predicate on $m$ signs and let $\mu\in[-1,1]^m$. The probability
that $P$ accepts independent random signs whose respective means are the coordinates of
$\mu$ equals the multilinear extension of $P$ evaluated at $\mu$.
\end{theorem}

\begin{definition}[Zero bounded cube function]
\lean{FABL.BoundedCubeFunction.zero}
\label{FABL.BoundedCubeFunction.zero}
\leanok
\uses{FABL.BoundedCubeFunction}
For each natural number $n$, the zero bounded cube function is the real-valued function
on the $n$-dimensional Boolean cube that takes the value $0$ at every point; its values
lie in $[-1,1]$.
\end{definition}

\begin{theorem}[Value of the zero function]
\lean{FABL.BoundedCubeFunction.zero_apply}
\label{FABL.BoundedCubeFunction.zero_apply}
\leanok
\uses{FABL.BoundedCubeFunction.zero}
\difficulty{1}
For a natural number $n$ and any point $x$ of the $n$-dimensional sign cube, the zero
bounded cube function has value $0$ at $x$.
\end{theorem}

\begin{theorem}[Multilinear extension at zero]
\lean{FABL.SignPredicate.multilinearExtension_zero}
\label{FABL.SignPredicate.multilinearExtension_zero}
\leanok
\uses{FABL.SignPredicate, FABL.SignPredicate.indicator, FABL.SignPredicate.multilinearExtension, FABL.fourierCoeff}
\difficulty{4}
For any Boolean predicate $P$ on finitely many signs, its multilinear extension at the
zero vector equals the Fourier coefficient of its indicator function at the empty set.
\end{theorem}

\begin{definition}[Predicate family baseline]
\lean{FABL.predicateFamilyBaseline}
\label{FABL.predicateFamilyBaseline}
\leanok
\uses{FABL.SignPredicate, FABL.SignPredicate.indicator, FABL.fourierCoeff}
Let $\mathcal{P}$ be a finite nonempty family of sign predicates. Its baseline is the
minimum, over predicates in $\mathcal{P}$, of the constant Fourier coefficient of the
predicate's acceptance indicator on its input cube.
\end{definition}

\begin{theorem}[Baseline bounded by the zero extension]
\lean{FABL.predicateFamilyBaseline_le_multilinearExtension_zero}
\label{FABL.predicateFamilyBaseline_le_multilinearExtension_zero}
\leanok
\uses{FABL.SignPredicate, FABL.SignPredicate.indicator, FABL.SignPredicate.multilinearExtension, FABL.SignPredicate.multilinearExtension_zero, FABL.fourierCoeff, FABL.predicateFamilyBaseline}
\difficulty{1}
Let $\mathcal P$ be a nonempty finite family of sign predicates, and choose
$P\in\mathcal P$. The minimum, over $Q\in\mathcal P$, of the empty-set Fourier
coefficient of its acceptance indicator is at most the multilinear extension of $P$ at
the zero vector.
\end{theorem}

\begin{theorem}[Baseline bounded by the zero-assignment expectation]
\lean{FABL.predicateFamilyBaseline_le_zeroAssignmentExpectation}
\label{FABL.predicateFamilyBaseline_le_zeroAssignmentExpectation}
\leanok
\uses{FABL.SignPredicate, FABL.SignPredicate.randomizedValue, FABL.SignPredicate.randomizedValue_eq_multilinearExtension, FABL.pmfExpectation, FABL.pmfExpectation_const, FABL.pmfExpectation_mono, FABL.predicateFamilyBaseline, FABL.predicateFamilyBaseline_le_multilinearExtension_zero}
\difficulty{4}
Let $\mathcal P$ be a nonempty finite family of sign predicates, let $\Omega$ be a
finite set equipped with a probability law, and let $a:\Omega\to\mathcal P$ select a
predicate for each outcome. The minimum empty-set Fourier coefficient among predicates
in $\mathcal P$ is at most the expectation, under this law, of the randomized value of
$a(\omega)$ when every input coordinate has mean zero.
\end{theorem}

\begin{theorem}[Fourier formula for randomized signs]
\lean{FABL.hastadAcceptanceProbability_eq_fourierSum_randomizedSigns}
\proofsource{Analysis-of-Boolean-Functions-by-Ryan-ODonnell}{
  pdf-ffbe9c7b3306-p194-b002,
  pdf-ffbe9c7b3306-p194-b003,
  pdf-ffbe9c7b3306-p194-b004
}
\label{FABL.hastadAcceptanceProbability_eq_fourierSum_randomizedSigns}
\leanok
\uses{FABL.BoundedCubeFunction, FABL.fourierCoeff, FABL.hastadAcceptanceProbability, FABL.hastadAcceptanceProbability_eq_fourierSum}
\difficulty{1}
Let $n$ be a natural number, let $0\leq\delta\leq1$, and let $f:\{-1,1\}^n\to[-1,1]$.
The Håstad tester acceptance probability with noise parameter $\delta$ equals
$\frac12+\frac12\sum_{S\subseteq[n],\,|S|\text{ odd}}(1-\delta)^{|S|}\widehat f(S)^3$,
where $\widehat f(S)=\mathbb E_x[f(x)\prod_{i\in S}x_i]$ for uniform $x\in\{-1,1\}^n$.
\end{theorem}

\begin{theorem}[Hastad full randomized assignment dictator vs no notables]
\lean{FABL.hastad_fullRandomizedAssignment_dictatorVsNoNotables}
\statementsource{Analysis-of-Boolean-Functions-by-Ryan-ODonnell}{pdf-ffbe9c7b3306-p199-b006}
\label{FABL.hastad_fullRandomizedAssignment_dictatorVsNoNotables}
\leanok
\uses{FABL.IsDictatorVsNoNotablesTest, FABL.hastadDictatorTestFamily, FABL.hastadSoundnessError, FABL.hastad_dictatorVsNoNotables, FABL.maxE3LinPredicates}
\difficulty{1}
% dataset-pipeline: stub -- prose pending build_informal_statements.py
Statement pending regeneration.
\end{theorem}

\begin{theorem}[Stable influence le influence]
\lean{FABL.stableInfluence_le_influence}
\statementsource{Analysis-of-Boolean-Functions-by-Ryan-ODonnell}{
  pdf-ffbe9c7b3306-p057-b001,
  pdf-ffbe9c7b3306-p057-b003,
  pdf-ffbe9c7b3306-p264-b003
}
\label{FABL.stableInfluence_le_influence}
\leanok
\uses{FABL.influence, FABL.influence_eq_sum_sq_fourierCoeff, FABL.stableInfluence}
\difficulty{4}
% dataset-pipeline: stub -- prose pending build_informal_statements.py
Statement pending regeneration.
\end{theorem}

\begin{theorem}[Has small stable influences]
\lean{FABL.HasSmallInfluences.hasSmallStableInfluences}
\statementsource{Analysis-of-Boolean-Functions-by-Ryan-ODonnell}{pdf-ffbe9c7b3306-p145-b008}
\label{FABL.HasSmallInfluences.hasSmallStableInfluences}
\leanok
\uses{FABL.HasSmallInfluences, FABL.HasSmallStableInfluences, FABL.influence, FABL.stableInfluence, FABL.stableInfluence_le_influence, FABL.stableInfluence_one_eq_influence}
\difficulty{4}
% dataset-pipeline: stub -- prose pending build_informal_statements.py
Statement pending regeneration.
\end{theorem}

\begin{definition}[Predicate with negated inputs]
\lean{FABL.SignPredicate.negateInputs}
\label{FABL.SignPredicate.negateInputs}
\leanok
\uses{FABL.SignPredicate}
Let $P$ be a sign predicate of arity $r$, and choose a Boolean mask $\sigma$ on its $r$
input coordinates. The predicate obtained by negating inputs accepts an input $y$
exactly when $P$ accepts the input whose $i$th coordinate is $-y_i$ if $\sigma(i)$ is
true, and $y_i$ otherwise.
\end{definition}

\begin{definition}[Closure under input negations]
\lean{FABL.IsClosedUnderInputNegations}
\statementsource{Analysis-of-Boolean-Functions-by-Ryan-ODonnell}{pdf-ffbe9c7b3306-p199-b007}
\label{FABL.IsClosedUnderInputNegations}
\leanok
\uses{FABL.SignPredicate, FABL.SignPredicate.negateInputs}
A set of sign predicates is closed under input negations if, for every predicate in the
set and every choice of which input positions to negate, the predicate obtained by
negating those inputs also belongs to the set.
\end{definition}

\begin{definition}[Negated query check]
\lean{FABL.SignQueryCheck.negateQueries}
\statementsource{Analysis-of-Boolean-Functions-by-Ryan-ODonnell}{pdf-ffbe9c7b3306-p200-b001}
\label{FABL.SignQueryCheck.negateQueries}
\leanok
\uses{FABL.SignPredicate.negateInputs, FABL.SignQueryCheck}
Given a sign query check and a choice of predicate input positions to negate, form a new
check using the predicate with those inputs negated. At each selected position, negate
the corresponding queried sign vector; at every other position, keep the original query.
\end{definition}

\begin{theorem}[Closure under negated queries]
\lean{FABL.SignQueryCheck.negateQueries_mem}
\label{FABL.SignQueryCheck.negateQueries_mem}
\leanok
\uses{FABL.IsClosedUnderInputNegations, FABL.SignPredicate, FABL.SignQueryCheck, FABL.SignQueryCheck.negateQueries}
\difficulty{1}
Let $\Psi$ be a set of sign predicates closed under independently negating any input,
and let $C$ be a sign query check whose predicate belongs to $\Psi$. For every choice of
which input positions to negate, the predicate of the resulting check also belongs to
$\Psi$.
\end{theorem}

\begin{definition}[Coordinatewise sign-flip bijection]
\lean{FABL.signInputNegationEquiv}
\label{FABL.signInputNegationEquiv}
\leanok
\uses{FABL.Sign}
For a natural number $r$ and a Boolean mask $\sigma$ on $r$ coordinates, the
coordinatewise sign-flip map on sign-valued $r$-tuples sends $y_i$ to $-y_i$ when
$\sigma(i)$ is true and leaves it unchanged otherwise; this map is a bijection.
\end{definition}

\begin{theorem}[Monomial under sign changes]
\lean{FABL.monomial_signInputNegationEquiv}
\label{FABL.monomial_signInputNegationEquiv}
\leanok
\uses{FABL.Sign, FABL.monomial, FABL.signInputNegationEquiv, FABL.signValue_neg}
\difficulty{4}
Let $r$ be a nonnegative integer, let $y\in\{\pm1\}^r$, let
$S\subseteq\{0,\ldots,r-1\}$, and choose a set of coordinates to negate, specified by a
map $\sigma:\{0,\ldots,r-1\}\to\{\mathrm{false},\mathrm{true}\}$. The monomial indexed
by $S$ after these sign changes equals the original monomial multiplied by $\prod_{i\in
S}(-1)$ over the coordinates for which $\sigma(i)=\mathrm{true}$, with factor $1$ for
the other coordinates.
\end{theorem}

\begin{theorem}[Fourier coefficients under input negation]
\lean{FABL.fourierCoeff_indicator_negateInputs}
\label{FABL.fourierCoeff_indicator_negateInputs}
\leanok
\uses{FABL.Sign, FABL.SignPredicate, FABL.SignPredicate.indicator, FABL.SignPredicate.negateInputs, FABL.fourierCoeff, FABL.monomial, FABL.monomial_signInputNegationEquiv, FABL.signInputNegationEquiv}
\difficulty{5}
Let $P$ be a Boolean predicate on $n$ sign-valued inputs, let $\sigma$ specify a choice
of whether to negate each input, and let $S$ be a subset of the input coordinates. The
Fourier coefficient at $S$ of the indicator of the predicate after these input negations
equals the Fourier coefficient at $S$ of the original predicate's indicator, multiplied
by $\prod_{i\in S}(-1)$ over coordinates where $\sigma_i$ prescribes negation and by $1$
over the other coordinates.
\end{theorem}

\begin{theorem}[Negating inputs in the multilinear extension]
\lean{FABL.SignPredicate.multilinearExtension_negateInputs}
\label{FABL.SignPredicate.multilinearExtension_negateInputs}
\leanok
\uses{FABL.SignPredicate, FABL.SignPredicate.indicator, FABL.SignPredicate.multilinearExtension, FABL.SignPredicate.negateInputs, FABL.fourierCoeff, FABL.fourierCoeff_indicator_negateInputs}
\difficulty{4}
Let $P$ be a Boolean predicate on $m$ signs, let $\sigma_i\in\{0,1\}$ specify which
inputs to negate, and let $\mu\in\mathbb R^m$. The multilinear extension of the
predicate obtained by negating the selected inputs, evaluated at $\mu$, equals the
multilinear extension of $P$ evaluated at the vector whose $i$th coordinate is $-\mu_i$
when $\sigma_i=1$ and $\mu_i$ otherwise.
\end{theorem}

\begin{theorem}[Negating inputs in randomized evaluation]
\lean{FABL.SignPredicate.randomizedValue_negateInputs}
\label{FABL.SignPredicate.randomizedValue_negateInputs}
\leanok
\uses{FABL.SignPredicate, FABL.SignPredicate.multilinearExtension_negateInputs, FABL.SignPredicate.negateInputs, FABL.SignPredicate.randomizedValue, FABL.SignPredicate.randomizedValue_eq_multilinearExtension}
\difficulty{3}
Let $P$ be a Boolean predicate on $m$ signs, choose which inputs to negate, and let
$\mu\in[-1,1]^m$. The randomized value of the predicate with those inputs negated at
means $\mu$ equals the randomized value of $P$ at the vector obtained by negating the
corresponding coordinates of $\mu$.
\end{theorem}

\begin{definition}[Odd part of a cube function]
\lean{FABL.oddPart}
\statementsource{Analysis-of-Boolean-Functions-by-Ryan-ODonnell}{
  pdf-ffbe9c7b3306-p015-b015,
  pdf-ffbe9c7b3306-p035-b008,
  pdf-ffbe9c7b3306-p200-b001
}
\label{FABL.oddPart}
\leanok
For a real-valued function $f$ on the sign cube $\\{-1,1\\}^n$, its odd part is the
function $x\\mapsto (f(x)-f(-x))/2$.
\end{definition}

\begin{definition}[Odd part of a bounded cube function]
\lean{FABL.BoundedCubeFunction.oddPart}
\statementsource{Analysis-of-Boolean-Functions-by-Ryan-ODonnell}{
  pdf-ffbe9c7b3306-p015-b015,
  pdf-ffbe9c7b3306-p035-b008,
  pdf-ffbe9c7b3306-p200-b001
}
\label{FABL.BoundedCubeFunction.oddPart}
\leanok
\uses{FABL.BoundedCubeFunction, FABL.oddPart}
For a bounded real-valued function $f$ on the sign cube $\{-1,1\}^n$, its odd part is
the bounded function $x\\mapsto (f(x)-f(-x))/2$.
\end{definition}

\begin{theorem}[Extensionality of bounded cube functions]
\lean{FABL.BoundedCubeFunction.ext}
\label{FABL.BoundedCubeFunction.ext}
\leanok
\uses{FABL.BoundedCubeFunction}
\difficulty{3}
For a natural number $n$, two real-valued functions on the $n$-dimensional sign cube
that agree at every input are equal as bounded cube functions.
\end{theorem}

\begin{theorem}[Pointwise formula for the odd part]
\lean{FABL.BoundedCubeFunction.oddPart_apply}
\proofsource{Analysis-of-Boolean-Functions-by-Ryan-ODonnell}{pdf-ffbe9c7b3306-p035-b008}
\label{FABL.BoundedCubeFunction.oddPart_apply}
\leanok
\uses{FABL.BoundedCubeFunction, FABL.BoundedCubeFunction.oddPart}
\difficulty{1}
Let $n$ be a natural number, let $f$ be a bounded real-valued function on the
$n$-dimensional sign cube, and let $x$ belong to that cube. The odd part of $f$ at $x$
is $(f(x)-f(-x))/2$.
\end{theorem}

\begin{theorem}[Oddness of the odd part]
\lean{FABL.odd_oddPart}
\proofsource{Analysis-of-Boolean-Functions-by-Ryan-ODonnell}{pdf-ffbe9c7b3306-p035-b009}
\label{FABL.odd_oddPart}
\leanok
\uses{FABL.oddPart}
\difficulty{3}
For every natural number $n$ and every real-valued function $f$ on the sign cube
$\{-1,1\}^n$, the function $x\mapsto \bigl(f(x)-f(-x)\bigr)/2$ is odd; that is, its
value at $-x$ is the negative of its value at $x$.
\end{theorem}

\begin{theorem}[Odd part of a coordinate function]
\lean{FABL.BoundedCubeFunction.oddPart_ofBoolean_dictator}
\label{FABL.BoundedCubeFunction.oddPart_ofBoolean_dictator}
\leanok
\uses{FABL.BoundedCubeFunction.ext, FABL.BoundedCubeFunction.oddPart, FABL.BoundedCubeFunction.ofBoolean, FABL.dictator, FABL.signValue, FABL.signValue_neg}
\difficulty{3}
For a natural number $n$ and a coordinate $i\in\{0,\ldots,n-1\}$, the odd part of
the bounded real-valued function induced by the $i$th coordinate function on the sign
cube equals that function.
\end{theorem}

\begin{definition}[Odd symmetrized mean]
\lean{FABL.oddSymmetrizedMean}
\statementsource{Analysis-of-Boolean-Functions-by-Ryan-ODonnell}{pdf-ffbe9c7b3306-p200-b001}
\label{FABL.oddSymmetrizedMean}
\leanok
\uses{FABL.BoundedCubeFunction}
Given a function $f$ on the sign cube with values in $[-1,1]$, a point $x$ in the cube,
and a Boolean choice, the odd symmetrized value is $-f(-x)$ when the choice is true and
$f(x)$ when it is false.
\end{definition}

\begin{theorem}[Bounds for the odd symmetrized value]
\lean{FABL.oddSymmetrizedMean_mem_Icc}
\label{FABL.oddSymmetrizedMean_mem_Icc}
\leanok
\uses{FABL.BoundedCubeFunction, FABL.oddSymmetrizedMean}
\difficulty{4}
Let $n$ be a nonnegative integer, let $f$ be a real-valued function on the sign cube of
dimension $n$ with $-1\le f(y)\le 1$ for every cube point $y$, and let $x$ be any point
of that cube. For either choice of whether to negate, the value is $f(x)$ when no
negation is chosen and $-f(-x)$ when negation is chosen; in both cases it lies in
$[-1,1]$.
\end{theorem}

\begin{theorem}[Product average for independent Boolean choices]
\lean{FABL.expect_boolChoices_prod}
\label{FABL.expect_boolChoices_prod}
\leanok
\uses{FABL.expect_prod_finArrow}
\difficulty{5}
Let $r$ be a nonnegative integer, let
$a_i:\{\mathrm{false},\mathrm{true}\}\to\mathbb{R}$ for each $i\in\{1,\ldots,r\}$, and
let $S\subseteq\{1,\ldots,r\}$. For a uniformly chosen Boolean vector
$\sigma\in\{\mathrm{false},\mathrm{true}\}^r$, the average of $\prod_{i\in
S}a_i(\sigma_i)$ equals $\prod_{i\in
S}\frac{a_i(\mathrm{false})+a_i(\mathrm{true})}{2}$.
\end{theorem}

\begin{theorem}[Boolean averaging for the multilinear extension]
\lean{FABL.expect_multilinearExtension_boolChoices}
\proofsource{Analysis-of-Boolean-Functions-by-Ryan-ODonnell}{
  pdf-ffbe9c7b3306-p035-b004,
  pdf-ffbe9c7b3306-p199-b003
}
\label{FABL.expect_multilinearExtension_boolChoices}
\leanok
\uses{FABL.SignPredicate, FABL.SignPredicate.indicator, FABL.SignPredicate.multilinearExtension, FABL.expect_boolChoices_prod, FABL.fourierCoeff}
\difficulty{4}
Let $P$ be a Boolean predicate on inputs with $m$ sign coordinates, and for each
coordinate $i$ let $a_i$ assign a real number to each of the two Boolean choices. The
uniform expectation over all Boolean choice vectors $\sigma$ of the multilinear
extension of $P$ evaluated at the coordinates $a_i(\sigma_i)$ equals that extension
evaluated at the coordinatewise averages $(a_i(0)+a_i(1))/2$.
\end{theorem}

\begin{definition}[Odd symmetrized check value]
\lean{FABL.SignQueryCheck.oddSymmetrizedValue}
\statementsource{Analysis-of-Boolean-Functions-by-Ryan-ODonnell}{pdf-ffbe9c7b3306-p200-b001}
\label{FABL.SignQueryCheck.oddSymmetrizedValue}
\leanok
\uses{FABL.BoundedCubeFunction, FABL.SignQueryCheck, FABL.SignQueryCheck.negateQueries, FABL.SignQueryCheck.randomizedValue}
Given a sign query check and a bounded real-valued function on the sign cube, take the
uniform average over all choices of predicate input positions to negate of the check's
randomized value after negating the selected queries. This average is the odd
symmetrized value of the check on the function.
\end{definition}

\begin{theorem}[Odd symmetrization identity]
\lean{FABL.SignQueryCheck.oddSymmetrizedValue_eq_oddPart}
\proofsource{Analysis-of-Boolean-Functions-by-Ryan-ODonnell}{pdf-ffbe9c7b3306-p200-b001}
\label{FABL.SignQueryCheck.oddSymmetrizedValue_eq_oddPart}
\leanok
\uses{FABL.BoundedCubeFunction, FABL.BoundedCubeFunction.oddPart, FABL.SignPredicate.negateInputs, FABL.SignPredicate.randomizedValue, FABL.SignPredicate.randomizedValue_eq_multilinearExtension, FABL.SignPredicate.randomizedValue_negateInputs, FABL.SignQueryCheck, FABL.SignQueryCheck.negateQueries, FABL.SignQueryCheck.oddSymmetrizedValue, FABL.SignQueryCheck.randomizedValue, FABL.expect_multilinearExtension_boolChoices, FABL.oddSymmetrizedMean, FABL.oddSymmetrizedMean_mem_Icc}
\difficulty{5}
Let $C$ be a sign query check and let $f$ be a real-valued function on the sign cube.
Averaging the randomized value of $C$ over all independent choices to negate its queries
equals the randomized value of $C$ on the odd part of $f$, the function $x\mapsto
(f(x)-f(-x))/2$.
\end{theorem}

\begin{definition}[Odd symmetrized acceptance probability]
\lean{FABL.SignTester.oddSymmetrizedAcceptanceProbability}
\statementsource{Analysis-of-Boolean-Functions-by-Ryan-ODonnell}{pdf-ffbe9c7b3306-p200-b001}
\label{FABL.SignTester.oddSymmetrizedAcceptanceProbability}
\leanok
\uses{FABL.BoundedCubeFunction, FABL.SignQueryCheck.oddSymmetrizedValue, FABL.SignTester, FABL.pmfExpectation}
Given a finite-seed sign tester and a bounded real-valued function on the sign cube, the
odd symmetrized acceptance probability is the expectation, under the tester's seed law,
of the odd symmetrized value of the check selected by each seed.
\end{definition}

\begin{theorem}[Odd-symmetrized acceptance equals odd-part acceptance]
\lean{FABL.SignTester.oddSymmetrizedAcceptanceProbability_eq_oddPart}
\proofsource{Analysis-of-Boolean-Functions-by-Ryan-ODonnell}{pdf-ffbe9c7b3306-p200-b001}
\label{FABL.SignTester.oddSymmetrizedAcceptanceProbability_eq_oddPart}
\leanok
\uses{FABL.BoundedCubeFunction, FABL.BoundedCubeFunction.oddPart, FABL.SignQueryCheck.oddSymmetrizedValue_eq_oddPart, FABL.SignTester, FABL.SignTester.acceptanceProbability, FABL.SignTester.oddSymmetrizedAcceptanceProbability, FABL.pmfExpectation}
\difficulty{3}
Let $n$ be a natural number, let $T$ be a randomized sign tester on the $n$-dimensional
Boolean cube, and let $f$ be a function on that cube with values in $[-1,1]$. The
acceptance probability of $f$ under the odd-symmetrized tester equals the acceptance
probability under $T$ of the odd part $x\mapsto(f(x)-f(-x))/2$.
\end{theorem}

\begin{theorem}[Fourier coefficients under global sign reversal]
\lean{FABL.fourierCoeff_comp_neg}
\label{FABL.fourierCoeff_comp_neg}
\leanok
\uses{FABL.fourierCoeff, FABL.monomial, FABL.monomial_neg_input, FABL.signCubeNegEquiv}
\difficulty{4}
For every natural number $n$, every real-valued function $f$ on the $n$-dimensional sign
cube, and every coordinate set $S\subseteq\{1,\ldots,n\}$, the Fourier coefficient at
$S$ of the function $x\mapsto f(-x)$ equals $(-1)^{|S|}$ times the Fourier coefficient
of $f$ at $S$.
\end{theorem}

\begin{theorem}[Fourier coefficients of the odd part]
\lean{FABL.fourierCoeff_oddPart}
\label{FABL.fourierCoeff_oddPart}
\leanok
\uses{FABL.fourierCoeff, FABL.fourierCoeff_comp_neg, FABL.monomial, FABL.oddPart}
\difficulty{4}
Let $n$ be a natural number, let $f$ be a real-valued function on the $n$-dimensional
sign cube, and let $S$ be a set of coordinates. The Fourier coefficient at $S$ of the
odd part $x\mapsto (f(x)-f(-x))/2$ is $\bigl(1-(-1)^{|S|}\bigr)\widehat f(S)/2$.
\end{theorem}

\begin{theorem}[Even-degree coefficients of the odd part]
\lean{FABL.fourierCoeff_oddPart_of_even}
\label{FABL.fourierCoeff_oddPart_of_even}
\leanok
\uses{FABL.fourierCoeff, FABL.fourierCoeff_oddPart, FABL.oddPart}
\difficulty{2}
Let $n$ be a nonnegative integer, let $f$ be a real-valued function on $\\{-1,1\\}^n$,
and let $S\\subseteq\\{1,\\ldots,n\\}$ have even cardinality. Define the odd part of $f$
by $f_{\\mathrm{odd}}(x)=(f(x)-f(-x))/2$. Its Fourier coefficient at $S$, computed with
respect to the uniform measure on the cube and the character $\\prod_{i\\in S}x_i$, is
zero.
\end{theorem}

\begin{theorem}[Odd-degree coefficients of the odd part]
\lean{FABL.fourierCoeff_oddPart_of_odd}
\proofsource{Analysis-of-Boolean-Functions-by-Ryan-ODonnell}{pdf-ffbe9c7b3306-p035-b009}
\label{FABL.fourierCoeff_oddPart_of_odd}
\leanok
\uses{FABL.fourierCoeff, FABL.fourierCoeff_oddPart, FABL.oddPart}
\difficulty{2}
Let $n$ be a nonnegative integer, let $f$ be a real-valued function on $\\{-1,1\\}^n$,
and let $S\\subseteq\\{1,\\ldots,n\\}$ have odd cardinality. Define
$f_{\\mathrm{odd}}(x)=(f(x)-f(-x))/2$. The Fourier coefficient of $f_{\\mathrm{odd}}$ at
$S$, with respect to the uniform measure on the cube and the character $\\prod_{i\\in
S}x_i$, equals the Fourier coefficient of $f$ at $S$.
\end{theorem}

\begin{theorem}[Stable influence of the odd part]
\lean{FABL.stableInfluence_oddPart_le}
\label{FABL.stableInfluence_oddPart_le}
\leanok
\uses{FABL.fourierCoeff, FABL.fourierCoeff_oddPart_of_even, FABL.fourierCoeff_oddPart_of_odd, FABL.oddPart, FABL.stableInfluence}
\difficulty{4}
Let $n$ be a natural number, let $f:\\{−1,1\\}^n\\to\\mathbb R$, and let $i$ be a
coordinate. Define the odd part of $f$ by $f_{\\mathrm{odd}}(x)=(f(x)-f(-x))/2$. For
every $\\rho\\in[0,1]$, the stable influence of $f_{\\mathrm{odd}}$ at parameter $\\rho$
and coordinate $i$ is at most the stable influence of $f$ there.
\end{theorem}

\begin{theorem}[Odd symmetrization preserves parameters for a property]
\lean{FABL.SignTester.oddSymmetrization_preserves_parameters_of_property}
\proofsource{Analysis-of-Boolean-Functions-by-Ryan-ODonnell}{
  pdf-ffbe9c7b3306-p200-b001,
  pdf-ffbe9c7b3306-p200-b002
}
\label{FABL.SignTester.oddSymmetrization_preserves_parameters_of_property}
\leanok
\uses{FABL.BoundedCubeFunction, FABL.BoundedCubeFunction.oddPart, FABL.BoundedCubeFunction.oddPart_ofBoolean_dictator, FABL.BoundedCubeFunction.ofBoolean, FABL.SignTester, FABL.SignTester.acceptanceProbability, FABL.SignTester.oddSymmetrizedAcceptanceProbability, FABL.SignTester.oddSymmetrizedAcceptanceProbability_eq_oddPart, FABL.dictator, FABL.odd_oddPart}
\difficulty{4}
Let $n$ be a natural number, let $T$ be a randomized sign tester on the $n$-dimensional
Boolean cube, and let $P$ be a property of $[-1,1]$-valued functions on that cube.
Suppose $P(f)$ implies $P$ for the odd part $x\mapsto(f(x)-f(-x))/2$ of $f$, every
coordinate dictator is accepted by $T$ with probability at least $\beta$, and every odd
function $f$ with property $P$ is accepted by $T$ with probability at most $B$. Then
every coordinate dictator is accepted by the odd-symmetrized tester with probability at
least $\beta$, and every function with property $P$ is accepted by the odd-symmetrized
tester with probability at most $B$.
\end{theorem}

\begin{theorem}[Odd symmetrization preserves tester parameters]
\lean{FABL.SignTester.oddSymmetrization_preserves_parameters}
\proofsource{Analysis-of-Boolean-Functions-by-Ryan-ODonnell}{
  pdf-ffbe9c7b3306-p200-b001,
  pdf-ffbe9c7b3306-p200-b002
}
\label{FABL.SignTester.oddSymmetrization_preserves_parameters}
\leanok
\uses{FABL.BoundedCubeFunction, FABL.BoundedCubeFunction.ofBoolean, FABL.HasSmallStableInfluences, FABL.SignTester, FABL.SignTester.acceptanceProbability, FABL.SignTester.oddSymmetrization_preserves_parameters_of_property, FABL.SignTester.oddSymmetrizedAcceptanceProbability, FABL.dictator, FABL.oddPart, FABL.stableInfluence, FABL.stableInfluence_oddPart_le}
\difficulty{4}
Let $n$ be a natural number, let $T$ be a randomized sign tester on the $n$-dimensional
Boolean cube, and let $\beta,\varepsilon,\delta, B\in\mathbb{R}$ with
$0\leq\delta\leq1$. Suppose every coordinate dictator is accepted by $T$ with
probability at least $\beta$, and every odd function $f$ valued in $[-1,1]$ whose stable
influences at parameter $1-\delta$ are all at most $\varepsilon$ is accepted with
probability at most $B$. Then every coordinate dictator is accepted by the
odd-symmetrized tester with probability at least $\beta$, and every $[-1,1]$-valued
function whose stable influences at parameter $1-\delta$ are all at most $\varepsilon$
is accepted by the odd-symmetrized tester with probability at most $B$.
\end{theorem}

\begin{definition}[Original BLR tester]
\lean{FABL.originalBLRTester}
\statementsource{Analysis-of-Boolean-Functions-by-Ryan-ODonnell}{
  pdf-ffbe9c7b3306-p032-b002,
  pdf-ffbe9c7b3306-p159-b003,
  pdf-ffbe9c7b3306-p193-b003
}
\label{FABL.originalBLRTester}
\leanok
\uses{FABL.SignTester, FABL.e3LinPredicate, FABL.scaledPointwiseProduct, FABL.uniformPMF}
For each natural number $n$, define a sign tester whose seeds are pairs $(x,y)$ of
points chosen independently and uniformly from the sign cube $\{-1,1\}^n$. Its check
uses the three-query linearity predicate with parameter $1$, querying $x$, $y$, and
their coordinatewise product scaled by $1$.
\end{definition}

\begin{theorem}[BLR acceptance and Fourier coefficients]
\lean{FABL.originalBLRTester_acceptanceProbability_eq_fourierSum}
\proofsource{Analysis-of-Boolean-Functions-by-Ryan-ODonnell}{pdf-ffbe9c7b3306-p032-b004}
\label{FABL.originalBLRTester_acceptanceProbability_eq_fourierSum}
\leanok
\uses{FABL.BoundedCubeFunction, FABL.Sign, FABL.SignPredicate.randomizedValue, FABL.SignQueryCheck.randomizedValue, FABL.SignTester.acceptanceProbability, FABL.e3LinPredicate, FABL.e3LinPredicate_randomizedValue, FABL.expect_hastadCore_eq_fourierSum, FABL.fourierCoeff, FABL.fourier_expansion, FABL.monomial, FABL.noiseOperator, FABL.noiseOperator_fourier_expansion, FABL.originalBLRTester, FABL.pmfExpectation, FABL.pmfExpectation_bind, FABL.pmfExpectation_map, FABL.pmfExpectation_uniformPMF_eq_expect, FABL.scaledPointwiseProduct, FABL.signValue, FABL.signValue_one, FABL.uniformPMF}
\difficulty{6}
Let $n$ be a nonnegative integer and let $f$ be a real-valued function on the sign cube
$\\{-1,1\\}^n$ taking values in $[-1,1]$. For each subset $S$ of the coordinates, let
$\\widehat f(S)$ be the uniform average of $f(x)\\prod_{i\\in S}x_i$ over the sign cube.
The acceptance probability of the original three-query BLR tester on $f$ is
$\\frac12+\\frac12\\sum_{S\\subseteq\\{0,\\ldots,n-1\\}}\\widehat f(S)^3$.
\end{theorem}

\begin{theorem}[Dictator acceptance by the BLR tester]
\lean{FABL.originalBLRTester_acceptanceProbability_dictator}
\proofsource{Analysis-of-Boolean-Functions-by-Ryan-ODonnell}{pdf-ffbe9c7b3306-p032-b004}
\label{FABL.originalBLRTester_acceptanceProbability_dictator}
\leanok
\uses{FABL.BooleanFunction.toReal, FABL.BoundedCubeFunction.ofBoolean, FABL.SignTester.acceptanceProbability, FABL.dictator, FABL.dictator_toReal_eq_monomial_singleton, FABL.fourierCoeff, FABL.fourierCoeff_monomial, FABL.monomial, FABL.originalBLRTester, FABL.originalBLRTester_acceptanceProbability_eq_fourierSum}
\difficulty{3}
For every nonnegative integer $n$ and every coordinate $i\in\{0,\ldots,n-1\}$, the
original three-query BLR tester accepts the sign-valued coordinate function $x\mapsto
x_i$ with probability $1$.
\end{theorem}

\begin{definition}[Noise operator on bounded cube functions]
\lean{FABL.BoundedCubeFunction.noise}
\statementsource{Analysis-of-Boolean-Functions-by-Ryan-ODonnell}{
  pdf-ffbe9c7b3306-p055-b002,
  pdf-ffbe9c7b3306-p194-b001,
  pdf-ffbe9c7b3306-p194-b003
}
\label{FABL.BoundedCubeFunction.noise}
\leanok
\uses{FABL.BoundedCubeFunction, FABL.noiseKernel, FABL.noiseOperator, FABL.noiseOperator_apply_eq_pmfExpectation, FABL.pmfExpectation_const, FABL.pmfExpectation_mono}
For a natural number $n$, a real parameter $\rho\in[-1,1]$, and a function
$f:\{-1,1\}^n\to[-1,1]$, define its noise transform at $x\in\{-1,1\}^n$ to be the
expected value of $f$ on the point obtained by applying the noise kernel with parameter
$\rho$ to $x$. This transform again takes values in $[-1,1]$.
\end{definition}

\begin{theorem}[Noise operator evaluation]
\lean{FABL.BoundedCubeFunction.noise_apply}
\label{FABL.BoundedCubeFunction.noise_apply}
\leanok
\uses{FABL.BoundedCubeFunction, FABL.BoundedCubeFunction.noise, FABL.noiseOperator}
\difficulty{1}
Let $n$ be a natural number, let $
ho\in[-1,1]$, let $f$ be a bounded real-valued
function on the $n$-dimensional sign cube, and let $x$ be a point of that cube. Applying
the noise operator with parameter $
ho$ to $f$ and evaluating at $x$ gives the
noise-kernel expectation of $f$ at $x$.
\end{theorem}

\begin{theorem}[Multilinear extension under product distributions]
\lean{FABL.pmfExpectation_multilinearExtension_independentProduct}
\proofsource{Analysis-of-Boolean-Functions-by-Ryan-ODonnell}{
  pdf-ffbe9c7b3306-p035-b004,
  pdf-ffbe9c7b3306-p199-b003
}
\label{FABL.pmfExpectation_multilinearExtension_independentProduct}
\leanok
\uses{FABL.SignPredicate, FABL.SignPredicate.multilinearExtension, FABL.independentProductPMF, FABL.pmfExpectation, FABL.pmfExpectation_const_mul, FABL.pmfExpectation_independentProduct_subsetProd, FABL.pmfExpectation_sum}
\difficulty{3}
Let $P$ be a Boolean predicate on $m$ sign inputs, let each coordinate space $A_i$ be
finite with a probability mass function $p_i$, and let $a_i:A_i\to\mathbb{R}$. Form a
random tuple $Y=(Y_i)_{i=1}^m$ with independent coordinates distributed according to
$p_i$. The expectation of the multilinear extension of the acceptance indicator of $P$,
evaluated at the vector $(a_i(Y_i))_{i=1}^m$, equals that extension evaluated at the
vector of coordinate expectations $(\mathbb{E}_{p_i}[a_i])_{i=1}^m$.
\end{theorem}

\begin{definition}[Noise-preprocessed check value]
\lean{FABL.SignQueryCheck.noisePreprocessedValue}
\statementsource{Analysis-of-Boolean-Functions-by-Ryan-ODonnell}{pdf-ffbe9c7b3306-p200-b003}
\label{FABL.SignQueryCheck.noisePreprocessedValue}
\leanok
\uses{FABL.BoundedCubeFunction, FABL.SignPredicate.randomizedValue, FABL.SignQueryCheck, FABL.independentProductPMF, FABL.noiseKernel, FABL.pmfExpectation}
Let $n$ be a natural number, let $C$ be a check consisting of a sign-valued predicate
and a cube query at each predicate input, let $\delta\in[0,1]$, and let
$f:\{-1,1\}^n\to[-1,1]$. Independently noise each query using the noise kernel with
parameter $1-\delta$, then take the expected randomized value of the predicate when its
inputs are the values of $f$ on the resulting noisy queries.
\end{definition}

\begin{theorem}[Noise preprocessing and the noise operator]
\lean{FABL.SignQueryCheck.noisePreprocessedValue_eq_noiseOperator}
\proofsource{Analysis-of-Boolean-Functions-by-Ryan-ODonnell}{pdf-ffbe9c7b3306-p200-b003}
\label{FABL.SignQueryCheck.noisePreprocessedValue_eq_noiseOperator}
\leanok
\uses{FABL.BoundedCubeFunction, FABL.BoundedCubeFunction.noise, FABL.SignPredicate.multilinearExtension, FABL.SignPredicate.randomizedValue, FABL.SignPredicate.randomizedValue_eq_multilinearExtension, FABL.SignQueryCheck, FABL.SignQueryCheck.noisePreprocessedValue, FABL.SignQueryCheck.randomizedValue, FABL.independentProductPMF, FABL.noiseKernel, FABL.noiseOperator, FABL.noiseOperator_apply_eq_pmfExpectation, FABL.pmfExpectation, FABL.pmfExpectation_multilinearExtension_independentProduct}
\difficulty{4}
Let $C$ be a sign query check, let $\delta\in[0,1]$, and let $f$ be a real-valued
function on the $n$-dimensional sign cube taking values in $[-1,1]$. The value of $C$
after independently applying noise with correlation $1-\delta$ to its queries and then
evaluating $f$ equals the randomized value of $C$ on the function obtained by applying
the noise operator with correlation $1-\delta$ to $f$.
\end{theorem}

\begin{definition}[Noise-preprocessed acceptance probability]
\lean{FABL.SignTester.noisePreprocessedAcceptanceProbability}
\statementsource{Analysis-of-Boolean-Functions-by-Ryan-ODonnell}{pdf-ffbe9c7b3306-p200-b003}
\label{FABL.SignTester.noisePreprocessedAcceptanceProbability}
\leanok
\uses{FABL.BoundedCubeFunction, FABL.SignQueryCheck.noisePreprocessedValue, FABL.SignTester, FABL.pmfExpectation}
Let $n$ be a natural number, let $T$ be a tester with a finite random seed, and let
$\delta\in[0,1]$. For a function $f:\{-1,1\}^n\to[-1,1]$, define the noise-preprocessed
acceptance probability as the expected noise-preprocessed check value, where the
expectation is over the tester's seed.
\end{definition}

\begin{theorem}[Noise preprocessing and tester acceptance]
\lean{FABL.SignTester.noisePreprocessedAcceptanceProbability_eq_noiseOperator}
\label{FABL.SignTester.noisePreprocessedAcceptanceProbability_eq_noiseOperator}
\leanok
\uses{FABL.BoundedCubeFunction, FABL.BoundedCubeFunction.noise, FABL.SignQueryCheck.noisePreprocessedValue_eq_noiseOperator, FABL.SignTester, FABL.SignTester.acceptanceProbability, FABL.SignTester.noisePreprocessedAcceptanceProbability, FABL.pmfExpectation}
\difficulty{3}
Let $T$ be a sign tester, let $\delta\in[0,1]$, and let $f$ be a real-valued function on
the sign cube taking values in $[-1,1]$. The acceptance probability of $T$ after
independently applying noise with correlation $1-\delta$ to each check's queries equals
the acceptance probability of $T$ on the function obtained by applying the noise
operator with correlation $1-\delta$ to $f$.
\end{theorem}

\begin{definition}[Noise-preprocessed tester]
\lean{FABL.SignTester.noisePreprocessed}
\statementsource{Analysis-of-Boolean-Functions-by-Ryan-ODonnell}{pdf-ffbe9c7b3306-p200-b003}
\label{FABL.SignTester.noisePreprocessed}
\leanok
\uses{FABL.SignTester, FABL.independentProductPMF, FABL.noiseKernel}
Let $n$ be a natural number, let $T$ be a tester with a finite seed and a probability
law on its seeds, and let $\delta\in[0,1]$. Define a new tester by first drawing a seed
according to the law of $T$, then independently applying noise with parameter $1-\delta$
to each query of the check selected by that seed; the resulting check uses the same
predicate and the noisy queries.
\end{definition}

\begin{theorem}[Acceptance probability after preprocessing]
\lean{FABL.SignTester.noisePreprocessed_acceptanceProbability}
\proofsource{Analysis-of-Boolean-Functions-by-Ryan-ODonnell}{pdf-ffbe9c7b3306-p200-b003}
\label{FABL.SignTester.noisePreprocessed_acceptanceProbability}
\leanok
\uses{FABL.BoundedCubeFunction, FABL.SignQueryCheck.randomizedValue, FABL.SignTester, FABL.SignTester.acceptanceProbability, FABL.SignTester.noisePreprocessed, FABL.SignTester.noisePreprocessedAcceptanceProbability, FABL.independentProductPMF, FABL.noiseKernel, FABL.pmfExpectation, FABL.pmfExpectation_bind, FABL.pmfExpectation_map}
\difficulty{4}
Let $T$ be a sign tester, let $\delta\in[0,1]$, and let $f$ be a real-valued function on
the sign cube taking values in $[-1,1]$. The acceptance probability of the tester formed
by preprocessing $T$ with independent query noise of correlation $1-\delta$ equals the
noise-preprocessed acceptance probability of $T$ on $f$.
\end{theorem}

\begin{theorem}[Predicate preservation under preprocessing]
\lean{FABL.SignTester.noisePreprocessed_usesPredicates}
\label{FABL.SignTester.noisePreprocessed_usesPredicates}
\leanok
\uses{FABL.SignPredicate, FABL.SignTester, FABL.SignTester.UsesPredicates, FABL.SignTester.noisePreprocessed}
\difficulty{1}
Let $T$ be a sign tester, let $\Psi$ be a set of sign predicates, and let
$\delta\in[0,1]$. If every check selected by $T$ has its predicate in $\Psi$, then every
check selected by the noise-preprocessed tester also has its predicate in $\Psi$.
\end{theorem}

\begin{theorem}[Odd-symmetrized BLR acceptance identity]
\lean{FABL.originalBLRTester_oddSymmetrizedAcceptanceProbability_eq_oddBLR}
\label{FABL.originalBLRTester_oddSymmetrizedAcceptanceProbability_eq_oddBLR}
\leanok
\uses{FABL.BoundedCubeFunction, FABL.SignTester.oddSymmetrizedAcceptanceProbability, FABL.SignTester.oddSymmetrizedAcceptanceProbability_eq_oddPart, FABL.fourierCoeff, FABL.fourierCoeff_oddPart_of_even, FABL.fourierCoeff_oddPart_of_odd, FABL.hastadAcceptanceProbability, FABL.hastadAcceptanceProbability_eq_fourierSum, FABL.oddBLRAcceptanceProbability, FABL.oddPart, FABL.originalBLRTester, FABL.originalBLRTester_acceptanceProbability_eq_fourierSum}
\difficulty{5}
For every nonnegative integer $n$ and every real-valued function $f$ on the sign cube
with values in $[-1,1]$, the acceptance probability of the original BLR tester after
independently symmetrizing each of its three queries by a uniform random sign flip
equals the acceptance probability of the zero-noise Håstad tester (the odd BLR tester)
on $f$.
\end{theorem}

\begin{theorem}[Noise commutes with taking the odd part]
\lean{FABL.BoundedCubeFunction.noise_oddPart}
\label{FABL.BoundedCubeFunction.noise_oddPart}
\leanok
\uses{FABL.BoundedCubeFunction, FABL.BoundedCubeFunction.ext, FABL.BoundedCubeFunction.noise, FABL.BoundedCubeFunction.oddPart, FABL.fourierCoeff_noiseOperator, FABL.fourierCoeff_oddPart, FABL.fourier_expansion, FABL.noiseOperator, FABL.oddPart}
\difficulty{4}
Let $n$ be a natural number, let $
ho\in[-1,1]$, and let $f$ be a bounded real-valued
function on the $n$-dimensional sign cube. The odd part of the noise of $f$ with
parameter $
ho$ equals the noise with parameter $
ho$ of the odd part of $f$.
\end{theorem}

\begin{definition}[Noise of a function with small stable influences]
\lean{FABL.IsNoiseOfSmallStableInfluences}
\statementsource{Analysis-of-Boolean-Functions-by-Ryan-ODonnell}{pdf-ffbe9c7b3306-p195-b001}
\label{FABL.IsNoiseOfSmallStableInfluences}
\leanok
\uses{FABL.BoundedCubeFunction, FABL.BoundedCubeFunction.noise, FABL.HasSmallStableInfluences}
Let $n$ be a natural number, let $\rho\in[-1,1]$, let $\varepsilon\in\mathbb R$, and let
$g:\{-1,1\}^n\to[-1,1]$. The function $g$ is the noise transform at parameter $\rho$ of
some function $f:\{-1,1\}^n\to[-1,1]$ whose stable influences are small with both
parameters equal to $\varepsilon$.
\end{definition}

\begin{theorem}[Odd parts preserve the noise representation]
\lean{FABL.IsNoiseOfSmallStableInfluences.oddPart}
\label{FABL.IsNoiseOfSmallStableInfluences.oddPart}
\leanok
\uses{FABL.BoundedCubeFunction, FABL.BoundedCubeFunction.noise_oddPart, FABL.BoundedCubeFunction.oddPart, FABL.IsNoiseOfSmallStableInfluences, FABL.oddPart, FABL.stableInfluence, FABL.stableInfluence_oddPart_le}
\difficulty{4}
Let $n$ be a nonnegative integer, let $\\rho\\in[-1,1]$, and let
$\\varepsilon\\in[0,1]$. Suppose $g:\\{-1,1\\}^n\\to\\mathbb R$ takes values in $[-1,1]$
and is the noise transform at correlation $\\rho$ of some function $f$ taking values in
$[-1,1]$ whose stable influence at parameter $1-\\varepsilon$ is at most $\\varepsilon$
at every coordinate. Then the odd part $x\\mapsto(g(x)-g(-x))/2$ is also the noise
transform at correlation $\\rho$ of a function taking values in $[-1,1]$ whose stable
influence at parameter $1-\\varepsilon$ is at most $\\varepsilon$ at every coordinate.
\end{theorem}

\begin{theorem}[Odd noisy-function acceptance bound]
\lean{FABL.originalBLRTester_acceptanceProbability_odd_noise_le}
\label{FABL.originalBLRTester_acceptanceProbability_odd_noise_le}
\leanok
\uses{FABL.BoundedCubeFunction, FABL.BoundedCubeFunction.noise, FABL.BoundedCubeFunction.oddPart, FABL.BoundedCubeFunction.sum_sq_fourierCoeff_le_one, FABL.HasSmallStableInfluences, FABL.SignTester.acceptanceProbability, FABL.fourierCoeff, FABL.fourierCoeff_noiseOperator, FABL.fourierCoeff_oddPart_of_even, FABL.fourierCoeff_oddPart_of_odd, FABL.noiseOperator, FABL.oddPart, FABL.originalBLRTester, FABL.originalBLRTester_acceptanceProbability_eq_fourierSum, FABL.stableInfluence}
\difficulty{6}
Let $n$ be a nonnegative integer, let $\\varepsilon\\in[0,1]$, let $\\rho\\in[-1,1]$
with $\\rho\\ge 0$, and suppose $\\varepsilon\\le 1-\\rho^2$. Let $f$ be a real-valued
function on the sign cube with values in $[-1,1]$. For each coordinate $i$, define its
stable influence at parameter $1-\\varepsilon$ by $\\sum_{S\\ni
i}(1-\\varepsilon)^{|S|-1}\\widehat f(S)^2$, where $\\widehat f(S)$ is the uniform
Fourier coefficient. If every such stable influence is at most $\\varepsilon$, then the
original three-query BLR tester accepts the odd part of the $\\rho$-noise of $f$ with
probability at most $\\frac12+\\frac{\\sqrt{\\varepsilon}}2$.
\end{theorem}

\begin{theorem}[Noise Stability Bound for Odd BLR]
\lean{FABL.oddBLRAcceptanceProbability_noise_le}
\proofsource{Analysis-of-Boolean-Functions-by-Ryan-ODonnell}{
  pdf-ffbe9c7b3306-p194-b001,
  pdf-ffbe9c7b3306-p194-b004,
  pdf-ffbe9c7b3306-p195-b001
}
\label{FABL.oddBLRAcceptanceProbability_noise_le}
\leanok
\uses{FABL.BoundedCubeFunction, FABL.BoundedCubeFunction.ext, FABL.BoundedCubeFunction.noise, FABL.BoundedCubeFunction.oddPart, FABL.HasSmallStableInfluences, FABL.IsNoiseOfSmallStableInfluences, FABL.IsNoiseOfSmallStableInfluences.oddPart, FABL.SignTester.oddSymmetrization_preserves_parameters_of_property, FABL.noiseOperator, FABL.oddBLRAcceptanceProbability, FABL.oddPart, FABL.originalBLRTester, FABL.originalBLRTester_acceptanceProbability_dictator, FABL.originalBLRTester_acceptanceProbability_odd_noise_le, FABL.originalBLRTester_oddSymmetrizedAcceptanceProbability_eq_oddBLR}
\difficulty{5}
Let $n$ be a natural number, let $\\varepsilon\\in[0,1]$, let $\\rho\\in[0,1]$, and
suppose $\\varepsilon\\le 1-\\rho^2$. For every function $f:\\{-1,1\\}^n\\to\\mathbb R$
taking values in $[-1,1]$, if for each coordinate $i$ the sum $\\sum_{S\\ni
i}(1-\\varepsilon)^{|S|-1}\\widehat f(S)^2$ is at most $\\varepsilon$, then the odd BLR
acceptance probability of the noise-smoothed function at noise parameter $\\rho$ is at
most $\\frac12+\\frac{\\sqrt{\\varepsilon}}2$.
\end{theorem}

\begin{theorem}[Probability of retaining a sign]
\lean{FABL.coordinateNoisePMF_apply_self_for_signPMF}
\proofsource{Analysis-of-Boolean-Functions-by-Ryan-ODonnell}{pdf-ffbe9c7b3306-p053-b003}
\label{FABL.coordinateNoisePMF_apply_self_for_signPMF}
\leanok
\uses{FABL.Sign, FABL.coordinateNoisePMF, FABL.correlationKeepProbability}
\difficulty{4}
Let $\rho\in[-1,1]$, and let $p=(1+\rho)/2$. For any sign $b\in\{−1,1\}$, the coordinate
noise distribution with parameter $\rho$, started at $b$, assigns probability $p$ to
$b$.
\end{theorem}

\begin{theorem}[Probability of flipping a sign]
\lean{FABL.coordinateNoisePMF_apply_neg_for_signPMF}
\proofsource{Analysis-of-Boolean-Functions-by-Ryan-ODonnell}{pdf-ffbe9c7b3306-p053-b003}
\label{FABL.coordinateNoisePMF_apply_neg_for_signPMF}
\leanok
\uses{FABL.Sign, FABL.coordinateNoisePMF, FABL.correlationKeepProbability}
\difficulty{4}
Let $\rho\in[-1,1]$, and let $p=(1+\rho)/2$. For any sign $b\in\{−1,1\}$, the coordinate
noise distribution with parameter $\rho$, started at $b$, assigns probability $1-p$ to
$-b$.
\end{theorem}

\begin{theorem}[Keep probability at one]
\lean{FABL.correlationKeepProbability_one_coe}
\label{FABL.correlationKeepProbability_one_coe}
\leanok
\uses{FABL.correlationKeepProbability}
\difficulty{3}
Given that $1\in[-1,1]$, the keep probability at correlation $\rho=1$, viewed as an
extended nonnegative real number, is one.
\end{theorem}

\begin{theorem}[Keep probability at negative one]
\lean{FABL.correlationKeepProbability_neg_one_coe}
\label{FABL.correlationKeepProbability_neg_one_coe}
\leanok
\uses{FABL.correlationKeepProbability}
\difficulty{3}
Given that $-1\in[-1,1]$, the keep probability at correlation $\rho=-1$, viewed as an
extended nonnegative real number, is zero.
\end{theorem}

\begin{theorem}[Probability of Retaining a Scaled Sign]
\lean{FABL.signPMFOfMean_scaled_sign_apply_self_toReal}
\label{FABL.signPMFOfMean_scaled_sign_apply_self_toReal}
\leanok
\uses{FABL.Sign, FABL.coordinateNoisePMF_apply_neg_for_signPMF, FABL.coordinateNoisePMF_apply_self_for_signPMF, FABL.correlationKeepProbability_le_one, FABL.signPMFOfMean, FABL.signValue, FABL.signValue_neg, FABL.signValue_one}
\difficulty{4}
Let $0\le\delta\le1$ and let $b\in\{-1,1\}$. Under the sign-valued noise distribution
whose mean is $(1-\delta)b$, the probability of obtaining $b$ is $1-\delta/2$.
\end{theorem}

\begin{theorem}[Sign distribution at a sign mean]
\lean{FABL.signPMFOfMean_signValue_eq_pure}
\label{FABL.signPMFOfMean_signValue_eq_pure}
\leanok
\uses{FABL.Sign, FABL.coordinateNoisePMF_apply_neg_for_signPMF, FABL.coordinateNoisePMF_apply_self_for_signPMF, FABL.correlationKeepProbability_neg_one_coe, FABL.correlationKeepProbability_one_coe, FABL.signPMFOfMean, FABL.signValue, FABL.signValue_neg_one, FABL.signValue_one}
\difficulty{4}
For every sign $b\in\{\pm1\}$, the probability distribution on signs whose mean
parameter is the real value of $b$ is the point mass at $b$.
\end{theorem}

\begin{theorem}[Randomized value at sign means]
\lean{FABL.SignPredicate.randomizedValue_signValues}
\proofsource{Analysis-of-Boolean-Functions-by-Ryan-ODonnell}{
  pdf-ffbe9c7b3306-p199-b003,
  pdf-ffbe9c7b3306-p199-b006
}
\label{FABL.SignPredicate.randomizedValue_signValues}
\leanok
\uses{FABL.Sign, FABL.SignPredicate, FABL.SignPredicate.randomizedValue, FABL.independentProductPMF, FABL.independentProductPMF_apply, FABL.pmfExpectation, FABL.signPMFOfMean_signValue_eq_pure, FABL.signValue}
\difficulty{4}
Let $P$ be a Boolean predicate on $m$ signs and fix a sign vector $x\in\{\pm1\}^m$. The
randomized value of $P$ when each input mean is the corresponding coordinate of $x$
equals $1$ if $P$ accepts $x$ and $0$ otherwise.
\end{theorem}

\begin{theorem}[Lower bound for randomized acceptance]
\lean{FABL.SignPredicate.randomizedValue_scaled_sign_ge}
\label{FABL.SignPredicate.randomizedValue_scaled_sign_ge}
\leanok
\uses{FABL.Sign, FABL.SignPredicate, FABL.SignPredicate.randomizedValue, FABL.SignPredicate.randomizedValue_mem_Icc, FABL.independentProductPMF, FABL.independentProductPMF_apply, FABL.pmfExpectation, FABL.signPMFOfMean, FABL.signPMFOfMean_scaled_sign_apply_self_toReal, FABL.signValue}
\difficulty{5}
Let $P$ be a Boolean predicate on $m$ signs, let $\delta\in[0,1]$, and fix a sign vector
$x\in\{\pm1\}^m$. Set $\mu_i=(1-\delta)x_i$. Then the indicator that $P$ accepts $x$,
minus $m\delta/2$, is at most the randomized value of $P$ at means $\mu$.
\end{theorem}

\begin{theorem}[Noise preprocessed value dictator ge]
\lean{FABL.SignQueryCheck.noisePreprocessedValue_dictator_ge}
\statementsource{Analysis-of-Boolean-Functions-by-Ryan-ODonnell}{pdf-ffbe9c7b3306-p194-b004}
\label{FABL.SignQueryCheck.noisePreprocessedValue_dictator_ge}
\leanok
\uses{FABL.BooleanFunction.toReal, FABL.BoundedCubeFunction.ofBoolean, FABL.Sign, FABL.SignPredicate.randomizedValue_eq_multilinearExtension, FABL.SignPredicate.randomizedValue_scaled_sign_ge, FABL.SignPredicate.randomizedValue_signValues, FABL.SignQueryCheck, FABL.SignQueryCheck.noisePreprocessedValue, FABL.SignQueryCheck.noisePreprocessedValue_eq_noiseOperator, FABL.SignQueryCheck.randomizedValue, FABL.dictator, FABL.dictator_toReal_eq_monomial_singleton, FABL.monomial, FABL.noiseOperator, FABL.noiseOperator_monomial_apply, FABL.signValue}
\difficulty{5}
% dataset-pipeline: stub -- prose pending build_informal_statements.py
Statement pending regeneration.
\end{theorem}

\begin{theorem}[Noise preprocessed acceptance probability dictator ge]
\lean{FABL.SignTester.noisePreprocessedAcceptanceProbability_dictator_ge}
\statementsource{Analysis-of-Boolean-Functions-by-Ryan-ODonnell}{pdf-ffbe9c7b3306-p194-b004}
\label{FABL.SignTester.noisePreprocessedAcceptanceProbability_dictator_ge}
\leanok
\uses{FABL.BoundedCubeFunction.ofBoolean, FABL.SignQueryCheck.noisePreprocessedValue_dictator_ge, FABL.SignQueryCheck.randomizedValue, FABL.SignTester, FABL.SignTester.acceptanceProbability, FABL.SignTester.noisePreprocessedAcceptanceProbability, FABL.dictator, FABL.pmfExpectation, FABL.pmfExpectation_add, FABL.pmfExpectation_const, FABL.pmfExpectation_mono}
\difficulty{4}
% dataset-pipeline: stub -- prose pending build_informal_statements.py
Statement pending regeneration.
\end{theorem}

\begin{theorem}[Noise preprocessing preserves parameters]
\lean{FABL.SignTester.noisePreprocessing_preserves_parameters}
\statementsource{Analysis-of-Boolean-Functions-by-Ryan-ODonnell}{
  pdf-ffbe9c7b3306-p200-b003,
  pdf-ffbe9c7b3306-p200-b004
}
\label{FABL.SignTester.noisePreprocessing_preserves_parameters}
\leanok
\uses{FABL.BoundedCubeFunction, FABL.BoundedCubeFunction.noise, FABL.BoundedCubeFunction.ofBoolean, FABL.HasSmallStableInfluences, FABL.SignTester, FABL.SignTester.acceptanceProbability, FABL.SignTester.noisePreprocessedAcceptanceProbability, FABL.SignTester.noisePreprocessedAcceptanceProbability_dictator_ge, FABL.SignTester.noisePreprocessedAcceptanceProbability_eq_noiseOperator, FABL.dictator}
\difficulty{4}
% dataset-pipeline: stub -- prose pending build_informal_statements.py
Statement pending regeneration.
\end{theorem}

\begin{definition}[Odd BLR query noise rate]
\lean{FABL.oddBLRQueryNoiseRate}
\statementsource{Analysis-of-Boolean-Functions-by-Ryan-ODonnell}{pdf-ffbe9c7b3306-p194-b001}
\label{FABL.oddBLRQueryNoiseRate}
\leanok
For each real $\delta$, define the odd BLR query noise rate to be $2\delta/3$.
\end{definition}

\begin{definition}[Odd BLR query noise cutoff]
\lean{FABL.oddBLRQueryNoiseCutoff}
\statementsource{Analysis-of-Boolean-Functions-by-Ryan-ODonnell}{
  pdf-ffbe9c7b3306-p194-b001,
  pdf-ffbe9c7b3306-p195-b001
}
\label{FABL.oddBLRQueryNoiseCutoff}
\leanok
\uses{FABL.oddBLRQueryNoiseRate}
For each real $\delta$, define the odd BLR query noise cutoff to be $1-(1-2\delta/3)^2$.
\end{definition}

\begin{definition}[Odd BLR noise soundness error]
\lean{FABL.oddBLRQueryNoiseSoundnessError}
\statementsource{Analysis-of-Boolean-Functions-by-Ryan-ODonnell}{pdf-ffbe9c7b3306-p195-b001}
\label{FABL.oddBLRQueryNoiseSoundnessError}
\leanok
\uses{FABL.hastadSoundnessError, FABL.oddBLRQueryNoiseCutoff}
For each real $\delta$, define the odd BLR query noise soundness error as the soundness
error at cutoff $1-(1-2\delta/3)^2$.
\end{definition}

\begin{theorem}[Noise rate lies in the unit interval]
\lean{FABL.oddBLRQueryNoiseRate_mem_Icc}
\label{FABL.oddBLRQueryNoiseRate_mem_Icc}
\leanok
\uses{FABL.oddBLRQueryNoiseRate}
\difficulty{3}
If $0<\delta<1/2$, then the noise rate $2\delta/3$ lies in the closed interval $[0,1]$.
\end{theorem}

\begin{theorem}[Noise cutoff lies in the unit interval]
\lean{FABL.oddBLRQueryNoiseCutoff_mem_Icc}
\label{FABL.oddBLRQueryNoiseCutoff_mem_Icc}
\leanok
\uses{FABL.oddBLRQueryNoiseCutoff, FABL.oddBLRQueryNoiseRate, FABL.oddBLRQueryNoiseRate_mem_Icc}
\difficulty{4}
If $0<\delta<1/2$, then the noise cutoff $1-(1-2\delta/3)^2$ lies in the closed interval
$[0,1]$.
\end{theorem}

\begin{theorem}[Positivity of the noise cutoff]
\lean{FABL.oddBLRQueryNoiseCutoff_pos}
\label{FABL.oddBLRQueryNoiseCutoff_pos}
\leanok
\uses{FABL.oddBLRQueryNoiseCutoff, FABL.oddBLRQueryNoiseRate, FABL.oddBLRQueryNoiseRate_mem_Icc}
\difficulty{4}
If $0<\delta<1/2$, then the noise cutoff $1-(1-2\delta/3)^2$ is positive.
\end{theorem}

\begin{definition}[Noisy odd BLR test family]
\lean{FABL.oddBLRQueryNoiseTestFamily}
\statementsource{Analysis-of-Boolean-Functions-by-Ryan-ODonnell}{
  pdf-ffbe9c7b3306-p194-b002,
  pdf-ffbe9c7b3306-p194-b003
}
\label{FABL.oddBLRQueryNoiseTestFamily}
\leanok
\uses{FABL.SignTester, FABL.SignTester.noisePreprocessed, FABL.oddBLRQueryNoiseRate, FABL.oddBLRQueryNoiseRate_mem_Icc, FABL.oddBLRTester}
For a noise parameter $\delta$ with $0<\delta<\tfrac12$ and a positive integer $n$,
define a sign tester on $n$ coordinates by applying noise preprocessing at rate
determined by $\delta$ to the odd BLR tester on $n$ coordinates.
\end{definition}

\begin{theorem}[Odd blr query noise parameter bounds]
\lean{FABL.oddBLRQueryNoise_parameter_bounds}
\statementsource{Analysis-of-Boolean-Functions-by-Ryan-ODonnell}{
  pdf-ffbe9c7b3306-p194-b004,
  pdf-ffbe9c7b3306-p195-b001
}
\label{FABL.oddBLRQueryNoise_parameter_bounds}
\leanok
\uses{FABL.BoundedCubeFunction, FABL.BoundedCubeFunction.noise, FABL.BoundedCubeFunction.ofBoolean, FABL.HasSmallStableInfluences, FABL.SignTester.acceptanceProbability, FABL.SignTester.noisePreprocessed_acceptanceProbability, FABL.SignTester.noisePreprocessing_preserves_parameters, FABL.dictator, FABL.hastadAcceptanceProbability, FABL.hastadAcceptanceProbability_dictator, FABL.oddBLRAcceptanceProbability, FABL.oddBLRAcceptanceProbability_noise_le, FABL.oddBLRQueryNoiseCutoff, FABL.oddBLRQueryNoiseRate, FABL.oddBLRQueryNoiseRate_mem_Icc, FABL.oddBLRQueryNoiseTestFamily, FABL.oddBLRTester}
\difficulty{5}
% dataset-pipeline: stub -- prose pending build_informal_statements.py
Statement pending regeneration.
\end{theorem}

\begin{theorem}[Hastad dictator vs no notables via odd blr and query noise]
\lean{FABL.hastad_dictatorVsNoNotables_via_oddBLR_and_queryNoise}
\statementsource{Analysis-of-Boolean-Functions-by-Ryan-ODonnell}{
  pdf-ffbe9c7b3306-p192-b006,
  pdf-ffbe9c7b3306-p194-b002,
  pdf-ffbe9c7b3306-p194-b003,
  pdf-ffbe9c7b3306-p194-b004,
  pdf-ffbe9c7b3306-p195-b001
}
\label{FABL.hastad_dictatorVsNoNotables_via_oddBLR_and_queryNoise}
\leanok
\uses{FABL.IsDictatorVsNoNotablesTest, FABL.SignTester.UsesPredicates, FABL.SignTester.acceptanceProbability_mem_Icc, FABL.SignTester.noisePreprocessed, FABL.SignTester.noisePreprocessed_usesPredicates, FABL.e3LinPredicate, FABL.hastadSoundnessError, FABL.hastadSoundnessError_mem_Icc, FABL.hastadSoundnessError_tendsto_zero, FABL.maxE3LinPredicates, FABL.maxE3LinPredicates_finite, FABL.oddBLRQueryNoiseCutoff, FABL.oddBLRQueryNoiseCutoff_mem_Icc, FABL.oddBLRQueryNoiseCutoff_pos, FABL.oddBLRQueryNoiseRate, FABL.oddBLRQueryNoiseRate_mem_Icc, FABL.oddBLRQueryNoiseSoundnessError, FABL.oddBLRQueryNoiseTestFamily, FABL.oddBLRQueryNoise_parameter_bounds, FABL.oddBLRTester}
\difficulty{5}
% dataset-pipeline: stub -- prose pending build_informal_statements.py
Statement pending regeneration.
\end{theorem}

\begin{definition}[Three-input sign predicate]
\lean{FABL.oxrPredicate}
\statementsource{Analysis-of-Boolean-Functions-by-Ryan-ODonnell}{
  pdf-ffbe9c7b3306-p034-b005,
  pdf-ffbe9c7b3306-p201-b003
}
\label{FABL.oxrPredicate}
\leanok
\uses{FABL.Sign, FABL.SignPredicate}
Let $P$ be the Boolean predicate on triples of signs $(y_0,y_1,y_2)\text{,}$ where each
$y_i$ is a unit in $\bbr$. The predicate accepts exactly when $y_0=-1$ or $y_1y_2=-1$.
\end{definition}

\begin{theorem}[Indicator of predicate acceptance]
\lean{FABL.oxrPredicate_indicator_eq}
\proofsource{Analysis-of-Boolean-Functions-by-Ryan-ODonnell}{pdf-ffbe9c7b3306-p034-b005}
\label{FABL.oxrPredicate_indicator_eq}
\leanok
\uses{FABL.Sign, FABL.oxrPredicate, FABL.signValue}
\difficulty{4}
For any triple of signs $y=(y_0,y_1,y_2)$, the indicator of whether the predicate
accepts $y$ is
\[
\mathbf{1}_{\{\text{the predicate accepts
}y\}}=\frac34-\frac{s(y_0)}4-\frac{s(y_1)s(y_2)}4-\frac{s(y_0)s(y_1)s(y_2)}4,
\]
where $s$ denotes the real value of a sign.
\end{theorem}

\begin{theorem}[Randomized value of the predicate]
\lean{FABL.oxrPredicate_randomizedValue}
\label{FABL.oxrPredicate_randomizedValue}
\leanok
\uses{FABL.Sign, FABL.SignPredicate.randomizedValue, FABL.independentProductPMF, FABL.oxrPredicate, FABL.oxrPredicate_indicator_eq, FABL.pmfExpectation, FABL.pmfExpectation_add, FABL.pmfExpectation_const, FABL.pmfExpectation_const_mul, FABL.pmfExpectation_independentProduct_subsetProd, FABL.pmfExpectation_signPMFOfMean, FABL.signPMFOfMean, FABL.signValue}
\difficulty{5}
Let $\mu_0,\mu_1,\mu_2\in[-1,1]$ be the means of three independent random signs supplied
to the predicate. Its randomized acceptance probability is
\[
\frac34-\frac{\mu_0}{4}-\frac{\mu_1\mu_2}{4}-\frac{\mu_0\mu_1\mu_2}{4}.
\]
\end{theorem}

\begin{definition}[Negated OXR predicates]
\lean{FABL.maxOXRPredicates}
\statementsource{Analysis-of-Boolean-Functions-by-Ryan-ODonnell}{pdf-ffbe9c7b3306-p201-b003}
\label{FABL.maxOXRPredicates}
\leanok
\uses{FABL.SignPredicate, FABL.SignPredicate.negateInputs, FABL.oxrPredicate}
Let $\mathcal{P}$ be the set of predicates obtained from the three-input OXR predicate
by independently negating its inputs according to a map
$\sigma:\{0,1,2\}\to\{\mathrm{false},\mathrm{true}\}$.
\end{definition}

\begin{theorem}[Finiteness of maximum OXR predicates]
\lean{FABL.maxOXRPredicates_finite}
\label{FABL.maxOXRPredicates_finite}
\leanok
\uses{FABL.SignPredicate.negateInputs, FABL.maxOXRPredicates, FABL.oxrPredicate}
\difficulty{1}
The family of maximum OXR predicates is finite.
\end{theorem}

\begin{definition}[OXR falsifying sign tuple]
\lean{FABL.oxrFalsifyingSign}
\statementsource{Analysis-of-Boolean-Functions-by-Ryan-ODonnell}{
  pdf-ffbe9c7b3306-p201-b003,
  pdf-ffbe9c7b3306-p201-b006
}
\label{FABL.oxrFalsifyingSign}
\leanok
\uses{FABL.Sign}
For a Boolean map $\sigma:\{0,1,2\}\to\{\mathrm{false},\mathrm{true}\}$ and a sign
$b\in\{\pm1\}$, define a sign tuple by starting with $(1,b,b)$ and negating coordinate
$i$ exactly when $\sigma(i)$ is true.
\end{definition}

\begin{theorem}[Expected acceptance of a random three-literal clause]
\lean{FABL.expect_oxrE3SatReplacementClauses}
\label{FABL.expect_oxrE3SatReplacementClauses}
\leanok
\uses{FABL.Sign, FABL.SignPredicate.negateInputs, FABL.e3SatSignClause, FABL.expect_sign, FABL.oxrFalsifyingSign, FABL.oxrPredicate}
\difficulty{5}
Let $\sigma$ assign a Boolean choice to each of three positions, and let
$v=(v_0,v_1,v_2)$ be a triple of signs in $\{\pm1\}$. Choose $b$ uniformly from
$\{\pm1\}$ and form the clause whose falsifying signs are $(1,b,b)$, negating the sign
at each position where $\sigma$ is true. The expected acceptance indicator of this
clause on $v$ is $\frac12$ plus half the acceptance indicator of the triple obtained by
negating the inputs according to $\sigma$ for the predicate that accepts $(x_0,x_1,x_2)$
exactly when $x_0\ne 1$ or $x_1\ne x_2$.
\end{theorem}

\begin{theorem}[Expected randomized value of a signed clause]
\lean{FABL.expect_oxrE3SatClause_randomizedValue}
\label{FABL.expect_oxrE3SatClause_randomizedValue}
\leanok
\uses{FABL.Sign, FABL.SignPredicate.negateInputs, FABL.SignPredicate.randomizedValue, FABL.e3SatSignClause, FABL.expect_oxrE3SatReplacementClauses, FABL.expect_sign, FABL.independentProductPMF, FABL.oxrFalsifyingSign, FABL.oxrPredicate, FABL.pmfExpectation, FABL.pmfExpectation_add, FABL.pmfExpectation_const, FABL.signPMFOfMean}
\difficulty{4}
Let $\sigma:\{1,2,3\}\to\{\mathrm{false},\mathrm{true}\}$ and let
$\mu_1,\mu_2,\mu_3\in[-1,1]$. For each sign $b\in\{−1,1\}$, form the three signs
obtained from $(1,b,b)$ by negating coordinate $i$ exactly when $\sigma_i$ is true, and
consider the three-literal clause that rejects precisely this sign triple. The uniform
average over $b$ of the clause's randomized value at coordinate means $\mu$ equals $1/2$
plus one half the randomized value at $\mu$ of the OXR predicate with its inputs negated
according to $\sigma$.
\end{theorem}

\begin{definition}[OXR sign tester]
\lean{FABL.OXRSignTester}
\statementsource{Analysis-of-Boolean-Functions-by-Ryan-ODonnell}{
  pdf-ffbe9c7b3306-p201-b003,
  pdf-ffbe9c7b3306-p201-b006
}
\label{FABL.OXRSignTester}
\leanok
An OXR sign tester on $n$ coordinates consists of a finite seed set with a probability
distribution, together with, for each seed, a Boolean choice of whether to negate each
of three predicate inputs and three queries in $\{\pm1\}^n$.
\end{definition}

\begin{definition}[Conversion to a sign tester]
\lean{FABL.OXRSignTester.toSignTester}
\statementsource{Analysis-of-Boolean-Functions-by-Ryan-ODonnell}{pdf-ffbe9c7b3306-p201-b003}
\label{FABL.OXRSignTester.toSignTester}
\leanok
\uses{FABL.OXRSignTester, FABL.SignPredicate.negateInputs, FABL.SignTester, FABL.oxrPredicate}
Every OXR sign tester on $n$ coordinates determines a sign tester with the same seed set
and seed distribution: at each seed, the check uses the OXR predicate with the
prescribed input negations and the three queries assigned to that seed.
\end{definition}

\begin{definition}[Conversion to an E3SAT tester]
\lean{FABL.OXRSignTester.toE3SatTester}
\statementsource{Analysis-of-Boolean-Functions-by-Ryan-ODonnell}{pdf-ffbe9c7b3306-p201-b003}
\label{FABL.OXRSignTester.toE3SatTester}
\leanok
\uses{FABL.OXRSignTester, FABL.Sign, FABL.SignTester, FABL.e3SatSignClause, FABL.oxrFalsifyingSign, FABL.uniformPMF}
Every OXR sign tester on $n$ coordinates determines a sign tester whose seeds are pairs
consisting of an original seed and a uniformly random sign. The seed distribution first
samples the original seed and then the sign; the check uses the E3SAT sign clause formed
from the falsifying sign tuple specified by the original seed's negations and the
sampled sign, together with that original seed's three queries.
\end{definition}

\begin{theorem}[Converted tester uses negated OXR predicates]
\lean{FABL.OXRSignTester.toSignTester_usesMaxOXR}
\label{FABL.OXRSignTester.toSignTester_usesMaxOXR}
\leanok
\uses{FABL.OXRSignTester, FABL.OXRSignTester.toSignTester, FABL.SignTester.UsesPredicates, FABL.maxOXRPredicates}
\difficulty{1}
Let $T$ be an OXR sign tester on $
$ coordinates. Every predicate used by its associated sign tester is obtained from the
OXR predicate by independently choosing, for each of its three inputs, whether to negate
that input.
\end{theorem}

\begin{theorem}[Converted tester uses three-literal clauses]
\lean{FABL.OXRSignTester.toE3SatTester_usesMaxE3Sat}
\proofsource{Analysis-of-Boolean-Functions-by-Ryan-ODonnell}{pdf-ffbe9c7b3306-p201-b003}
\label{FABL.OXRSignTester.toE3SatTester_usesMaxE3Sat}
\leanok
\uses{FABL.OXRSignTester, FABL.OXRSignTester.toE3SatTester, FABL.SignTester.UsesPredicates, FABL.maxE3SatSignPredicates, FABL.oxrFalsifyingSign}
\difficulty{1}
Let $T$ be an OXR sign tester on $
$ coordinates. Every predicate used by its associated E3SAT sign tester is a three-input
clause that rejects exactly one specified triple of signs and accepts every other
triple.
\end{theorem}

\begin{theorem}[Acceptance probability to e3 sat tester]
\lean{FABL.OXRSignTester.acceptanceProbability_toE3SatTester}
\statementsource{Analysis-of-Boolean-Functions-by-Ryan-ODonnell}{pdf-ffbe9c7b3306-p201-b003}
\label{FABL.OXRSignTester.acceptanceProbability_toE3SatTester}
\leanok
\uses{FABL.BoundedCubeFunction, FABL.OXRSignTester, FABL.OXRSignTester.toE3SatTester, FABL.OXRSignTester.toSignTester, FABL.Sign, FABL.SignPredicate.negateInputs, FABL.SignPredicate.randomizedValue, FABL.SignQueryCheck.randomizedValue, FABL.SignTester.acceptanceProbability, FABL.e3SatSignClause, FABL.expect_oxrE3SatClause_randomizedValue, FABL.oxrFalsifyingSign, FABL.oxrPredicate, FABL.pmfExpectation, FABL.pmfExpectation_add, FABL.pmfExpectation_bind, FABL.pmfExpectation_const, FABL.pmfExpectation_const_mul, FABL.pmfExpectation_map, FABL.pmfExpectation_uniformPMF_eq_expect, FABL.uniformPMF}
\difficulty{5}
% dataset-pipeline: stub -- prose pending build_informal_statements.py
Statement pending regeneration.
\end{theorem}

\begin{definition}[OXR replacement error]
\lean{FABL.oxrE3SatReplacementError}
\statementsource{Analysis-of-Boolean-Functions-by-Ryan-ODonnell}{pdf-ffbe9c7b3306-p201-b003}
\label{FABL.oxrE3SatReplacementError}
\leanok
For a function $e:\mathbb{R}\to\mathbb{R}$ and $\varepsilon\in\mathbb{R}$, define the
replacement error to be $e(\varepsilon)/2$.
\end{definition}

\begin{theorem}[Affine conversion to three-literal clauses]
\lean{FABL.oxrDictatorTest_to_e3Sat_affine}
\proofsource{Analysis-of-Boolean-Functions-by-Ryan-ODonnell}{pdf-ffbe9c7b3306-p201-b003}
\label{FABL.oxrDictatorTest_to_e3Sat_affine}
\leanok
\uses{FABL.IsDictatorVsNoNotablesTest, FABL.OXRSignTester, FABL.OXRSignTester.acceptanceProbability_toE3SatTester, FABL.OXRSignTester.toE3SatTester, FABL.OXRSignTester.toE3SatTester_usesMaxE3Sat, FABL.OXRSignTester.toSignTester, FABL.maxE3SatSignPredicates, FABL.maxE3SatSignPredicates_finite, FABL.maxOXRPredicates, FABL.oxrE3SatReplacementError}
\difficulty{5}
Let $\alpha,\beta\in\mathbb R$, let $\lambda:\mathbb R\to\mathbb R$, and suppose a
family of finite-seed OXR testers, one for each positive dimension, has a finite allowed
set of predicates; satisfies $0<\alpha<\beta\le1$; and has an error function $\lambda$
taking $[0,1]$ into $[0,1]$ and tending to zero as its positive argument tends to zero.
Suppose every test uses an allowed predicate, every coordinate dictator is accepted with
probability at least $\beta$, and every function whose stable influences at correlation
$1-\varepsilon$ are all at most $\varepsilon$ is accepted with probability at most
$\alpha+\lambda(\varepsilon)$ for each $\varepsilon\in[0,1]$. The associated testers
obtained by replacing each OXR test with a randomized three-literal clause test have a
finite allowed predicate set, completeness threshold $\frac12+\frac\beta2$, soundness
threshold $\frac12+\frac\alpha2$, and error function
$\varepsilon\mapsto\lambda(\varepsilon)/2$, with the same predicate-use,
dictator-completeness, and low-influence soundness properties.
\end{theorem}

\begin{definition}[Odd parity sign triple]
\lean{FABL.oxrParityTriple}
\label{FABL.oxrParityTriple}
\leanok
\uses{FABL.Sign}
For signs $a,b\in\{\pm1\}$, the associated triple is $(a,b,-ab)$.
\end{definition}

\begin{theorem}[First coordinate of the parity triple]
\lean{FABL.oxrParityTriple_zero}
\label{FABL.oxrParityTriple_zero}
\leanok
\uses{FABL.Sign, FABL.oxrParityTriple}
\difficulty{1}
For any pair of signs $(a,b)$, the first coordinate of the associated parity triple is
$a$.
\end{theorem}

\begin{theorem}[Second coordinate of the parity triple]
\lean{FABL.oxrParityTriple_one}
\label{FABL.oxrParityTriple_one}
\leanok
\uses{FABL.Sign, FABL.oxrParityTriple}
\difficulty{1}
For any pair of signs $(a,b)$, the second coordinate of the triple $(a,b,-ab)$ is $b$.
\end{theorem}

\begin{theorem}[Third coordinate of the parity triple]
\lean{FABL.oxrParityTriple_two}
\label{FABL.oxrParityTriple_two}
\leanok
\uses{FABL.Sign, FABL.oxrParityTriple}
\difficulty{1}
For any pair of signs $(a,b)$, the third coordinate of the associated parity triple is
$-ab$.
\end{theorem}

\begin{definition}[Sign pair triple]
\lean{FABL.oxrPairTriple}
\label{FABL.oxrPairTriple}
\leanok
\uses{FABL.Sign}
For signs $a,b\in\{\pm1\}$, the associated triple is $(a,b,-b)$.
\end{definition}

\begin{theorem}[First coordinate of the pair triple]
\lean{FABL.oxrPairTriple_zero}
\label{FABL.oxrPairTriple_zero}
\leanok
\uses{FABL.Sign, FABL.oxrPairTriple}
\difficulty{1}
For any pair of signs $(a,b)$, the first coordinate of the triple $(a,b,-b)$ is $a$.
\end{theorem}

\begin{theorem}[Second coordinate of the pair triple]
\lean{FABL.oxrPairTriple_one}
\label{FABL.oxrPairTriple_one}
\leanok
\uses{FABL.Sign, FABL.oxrPairTriple}
\difficulty{1}
For any pair of signs $(a,b)$, the second coordinate of the triple $(a,b,-b)$ is $b$.
\end{theorem}

\begin{theorem}[Third coordinate of the pair triple]
\lean{FABL.oxrPairTriple_two}
\label{FABL.oxrPairTriple_two}
\leanok
\uses{FABL.Sign, FABL.oxrPairTriple}
\difficulty{1}
For any pair of signs $(a,b)$, the third coordinate of the triple $(a,b,-b)$ is $-b$.
\end{theorem}

\begin{definition}[Parity triple distribution]
\lean{FABL.oxrParityTriplePMF}
\label{FABL.oxrParityTriplePMF}
\leanok
\uses{FABL.Sign, FABL.oxrParityTriple, FABL.uniformPMF}
Choose a pair of signs uniformly and map it to the triple $(a,b,-ab)$. The resulting
probability mass function on sign-valued triples is the parity triple distribution.
\end{definition}

\begin{definition}[Pair-based OXR triple distribution]
\lean{FABL.oxrPairTriplePMF}
\label{FABL.oxrPairTriplePMF}
\leanok
\uses{FABL.Sign, FABL.oxrPairTriple, FABL.uniformPMF}
Let $S$ be the set of signs. The pair-based OXR triple distribution on triples of signs
is the law obtained by sampling a pair uniformly from $S\times S$ and applying the OXR
pair-to-triple map.
\end{definition}

\begin{definition}[Coordinate OXR triple distribution]
\lean{FABL.oxrCoordinateTriplePMF}
\label{FABL.oxrCoordinateTriplePMF}
\leanok
\uses{FABL.Sign, FABL.oxrPairTriplePMF, FABL.oxrParityTriplePMF, FABL.pBiasedBitPMF}
Let $S$ be the set of signs and let $\delta\in[0,1]$. The coordinate OXR triple
distribution on triples of signs is obtained by first sampling a Boolean choice with
bias parameter $\delta$, then sampling from the pair-based OXR triple distribution if
the choice is true and from the parity triple distribution otherwise.
\end{definition}

\begin{theorem}[Expectation of a biased bit]
\lean{FABL.pmfExpectation_pBiasedBitPMF_affine}
\label{FABL.pmfExpectation_pBiasedBitPMF_affine}
\leanok
\uses{FABL.pBiasedBitPMF, FABL.pmfExpectation, FABL.pmfExpectation_add, FABL.pmfExpectation_const, FABL.pmfExpectation_const_mul, FABL.pmfExpectation_pBiasedBitPMF_zeroOneValue, FABL.zeroOneValue}
\difficulty{4}
Let $0\leq \delta\leq 1$ and let $g$ be a real-valued function on the two-point set
$\{\mathrm{false},\mathrm{true}\}$. Under the probability law assigning mass $1-\delta$
to $\mathrm{false}$ and mass $\delta$ to $\mathrm{true}$, the expectation of $g$ is
$(1-\delta)g(\mathrm{false})+\delta g(\mathrm{true})$.
\end{theorem}

\begin{definition}[Selected sign value]
\lean{FABL.selectedSign}
\label{FABL.selectedSign}
\leanok
\uses{FABL.Sign, FABL.signValue}
For a Boolean choice $b$ and a sign $s$, define its selected value to be the value of
$s$ when $b$ is true, and $1$ when $b$ is false.
\end{definition}

\begin{theorem}[Selected-sign expectation for the parity triple]
\lean{FABL.pmfExpectation_oxrParityTriplePMF_selectedSign}
\label{FABL.pmfExpectation_oxrParityTriplePMF_selectedSign}
\leanok
\uses{FABL.expect_sign, FABL.oxrParityTriplePMF, FABL.pmfExpectation, FABL.pmfExpectation_map, FABL.pmfExpectation_uniformPMF_eq_expect, FABL.selectedSign, FABL.signValue}
\difficulty{4}
Let $(U,V,-UV)$ be a triple formed from two independent uniform signs. For Boolean
selectors $a,b,c$, replace each coordinate whose selector is false by $1$ and take the
product of the resulting three real signs. Its expectation is $(-1)^a$ if $a=b=c$, and
is $0$ otherwise.
\end{theorem}

\begin{theorem}[Selected-sign expectation for the pair triple]
\lean{FABL.pmfExpectation_oxrPairTriplePMF_selectedSign}
\label{FABL.pmfExpectation_oxrPairTriplePMF_selectedSign}
\leanok
\uses{FABL.expect_sign, FABL.oxrPairTriplePMF, FABL.pmfExpectation, FABL.pmfExpectation_map, FABL.pmfExpectation_uniformPMF_eq_expect, FABL.selectedSign, FABL.signValue}
\difficulty{4}
Let $(U,V,-V)$ be a triple formed from two independent uniform signs. For Boolean
selectors $a,b,c$, replace each coordinate whose selector is false by $1$ and take the
product of the resulting three real signs. Its expectation is $(-1)^b$ if $a$ is false
and $b=c$, and is $0$ otherwise.
\end{theorem}

\begin{theorem}[Selected-sign expectation for a mixed triple]
\lean{FABL.pmfExpectation_oxrCoordinateTriplePMF_selectedSign}
\label{FABL.pmfExpectation_oxrCoordinateTriplePMF_selectedSign}
\leanok
\uses{FABL.oxrCoordinateTriplePMF, FABL.oxrPairTriplePMF, FABL.oxrParityTriplePMF, FABL.pmfExpectation, FABL.pmfExpectation_bind, FABL.pmfExpectation_oxrPairTriplePMF_selectedSign, FABL.pmfExpectation_oxrParityTriplePMF_selectedSign, FABL.pmfExpectation_pBiasedBitPMF_affine, FABL.selectedSign}
\difficulty{4}
Let $\delta\in[0,1]$, and choose a sign triple by first choosing the pair construction
with probability $\delta$ and the parity construction otherwise; in the pair
construction the triple is $(U,V,-V)$, and in the parity construction it is $(U,V,-UV)$,
where $U,V$ are independent uniform signs. For Boolean selectors $a,b,c$, form the
product of the selected signs, replacing each unselected coordinate by $1$. Its
expectation is $0$ if $b\ne c$; if $b=c$, it is $-(1-\delta)$ when
$a=b=c=\mathrm{true}$, $0$ when $a=\mathrm{true}$ and $b=c=\mathrm{false}$, $-\delta$
when $a=\mathrm{false}$ and $b=c=\mathrm{true}$, and $1$ when all three selectors are
false.
\end{theorem}

\begin{definition}[OXR triple cube distribution]
\lean{FABL.oxrTripleCubePMF}
\label{FABL.oxrTripleCubePMF}
\leanok
\uses{FABL.Sign, FABL.independentProductPMF, FABL.oxrCoordinateTriplePMF}
Let $n$ be a natural number, let $\delta\in[0,1]$, and let $S$ be the set of signs. The
OXR triple cube distribution on arrays of signs indexed by
$\mathrm{Fin}(n)\times\mathrm{Fin}(3)$ is the independent product of $n$ copies of the
coordinate OXR triple distribution with parameter $\delta$.
\end{definition}

\begin{definition}[Query point from a sample]
\lean{FABL.oxrQueryPoint}
\statementsource{Analysis-of-Boolean-Functions-by-Ryan-ODonnell}{pdf-ffbe9c7b3306-p201-b003}
\label{FABL.oxrQueryPoint}
\leanok
\uses{FABL.Sign}
For a natural number $n$, let a sample assign a sign in $\{\pm1\}$ to each pair $(i,q)$
with $i\in\{0,\ldots,n-1\}$ and $q\in\{0,1,2\}$. For a query $q\in\{0,1,2\}$, the
associated point of the $n$-dimensional sign cube has coordinate $i$ equal to the
sample's sign at $(i,q)$.
\end{definition}

\begin{definition}[Odd sign tester]
\lean{FABL.oxrOddSignTester}
\statementsource{Analysis-of-Boolean-Functions-by-Ryan-ODonnell}{
  pdf-ffbe9c7b3306-p201-b003,
  pdf-ffbe9c7b3306-p201-b006
}
\label{FABL.oxrOddSignTester}
\leanok
\uses{FABL.OXRSignTester, FABL.Sign, FABL.oxrQueryPoint, FABL.oxrTripleCubePMF, FABL.uniformPMF}
Given a natural number $n$ and a real number $\delta\in[0,1]$, define a tester with
seeds consisting of a sample from the correlated triple-cube distribution with parameter
$\delta$ and an independent uniformly random Boolean sign choice for each of the three
queries. The seed's Boolean choices specify the predicate-input negations, and its three
queries are the corresponding sample query points, negated wherever the associated
Boolean choice is true.
\end{definition}

\begin{theorem}[Product for two nested subsets]
\lean{FABL.prod_piecewise_two_sets}
\label{FABL.prod_piecewise_two_sets}
\leanok
\difficulty{4}
Let $\iota$ be a finite type, let $A,B\subseteq\iota$ with $A\subseteq B$, and let $a,b$
be real numbers. The product over $i\in\iota$ of $a$ for $i\in A$, $b$ for $i\in
B\setminus A$, and $1$ otherwise is $a^{|A|}b^{|B|-|A|}$.
\end{theorem}

\begin{theorem}[Pmf expectation oxr triple cube pmf monomials]
\lean{FABL.pmfExpectation_oxrTripleCubePMF_monomials}
\statementsource{Analysis-of-Boolean-Functions-by-Ryan-ODonnell}{pdf-ffbe9c7b3306-p194-b004}
\label{FABL.pmfExpectation_oxrTripleCubePMF_monomials}
\leanok
\uses{FABL.Sign, FABL.monomial, FABL.oxrCoordinateTriplePMF, FABL.oxrQueryPoint, FABL.oxrTripleCubePMF, FABL.pmfExpectation, FABL.pmfExpectation_independentProductPMF_prod, FABL.pmfExpectation_oxrCoordinateTriplePMF_selectedSign, FABL.prod_piecewise_two_sets, FABL.selectedSign, FABL.signValue}
\difficulty{6}
% dataset-pipeline: stub -- prose pending build_informal_statements.py
Statement pending regeneration.
\end{theorem}

\begin{theorem}[First-query expectation and the constant Fourier coefficient]
\lean{FABL.pmfExpectation_oxrTripleCubePMF_first}
\label{FABL.pmfExpectation_oxrTripleCubePMF_first}
\leanok
\uses{FABL.fourierCoeff, FABL.fourier_expansion, FABL.monomial, FABL.oxrQueryPoint, FABL.oxrTripleCubePMF, FABL.pmfExpectation, FABL.pmfExpectation_const_mul, FABL.pmfExpectation_oxrTripleCubePMF_monomials, FABL.pmfExpectation_sum}
\difficulty{4}
Let $n$ be a natural number, let $\delta\in[0,1]$, and let $f:\{-1,1\}^n\to\mathbb R$.
Choose independently at each coordinate a sign triple by selecting the pair construction
with probability $\delta$ and the parity construction otherwise, where these
constructions are $(U,V,-V)$ and $(U,V,-UV)$ for independent uniform signs $U,V$. The
expectation of $f$ evaluated on the first query vector is the uniform average of $f$ on
$\{-1,1\}^n$, equivalently its Fourier coefficient at the empty set.
\end{theorem}

\begin{theorem}[Second and third query correlation]
\lean{FABL.pmfExpectation_oxrTripleCubePMF_second_mul_third}
\label{FABL.pmfExpectation_oxrTripleCubePMF_second_mul_third}
\leanok
\uses{FABL.fourierCoeff, FABL.fourier_expansion, FABL.monomial, FABL.noiseStability, FABL.noiseStability_eq_sum_rho_pow_mul_sq_fourierCoeff, FABL.oxrQueryPoint, FABL.oxrTripleCubePMF, FABL.pmfExpectation, FABL.pmfExpectation_const_mul, FABL.pmfExpectation_oxrTripleCubePMF_monomials, FABL.pmfExpectation_sum}
\difficulty{4}
Let $n$ be a natural number, let $\delta\in[0,1]$, and let $f:\{-1,1\}^n\to\mathbb R$.
Choose independently at each coordinate a sign triple by selecting $(U,V,-V)$ with
probability $\delta$ and $(U,V,-UV)$ otherwise, with independent uniform signs $U,V$.
The expectation of the product of $f$ on the second and third query vectors equals the
noise stability of $f$ at correlation $-\delta$.
\end{theorem}

\begin{theorem}[Pmf expectation oxr triple cube pmf triple product fourier]
\lean{FABL.pmfExpectation_oxrTripleCubePMF_tripleProduct_fourier}
\statementsource{Analysis-of-Boolean-Functions-by-Ryan-ODonnell}{pdf-ffbe9c7b3306-p194-b004}
\label{FABL.pmfExpectation_oxrTripleCubePMF_tripleProduct_fourier}
\leanok
\uses{FABL.fourierCoeff, FABL.fourier_expansion, FABL.monomial, FABL.oxrQueryPoint, FABL.oxrTripleCubePMF, FABL.pmfExpectation, FABL.pmfExpectation_const_mul, FABL.pmfExpectation_oxrTripleCubePMF_monomials, FABL.pmfExpectation_sum}
\difficulty{6}
% dataset-pipeline: stub -- prose pending build_informal_statements.py
Statement pending regeneration.
\end{theorem}

\begin{theorem}[Subset intersection expectation formula]
\lean{FABL.expectDeltaRandomSubset_intersection_eq_sum}
\label{FABL.expectDeltaRandomSubset_intersection_eq_sum}
\leanok
\uses{FABL.deltaRandomSubsetWeight, FABL.expectDeltaRandomSubset, FABL.sum_deltaRandomSubsetWeight_inter_eq}
\difficulty{4}
Let $n$ be a nonnegative integer, let $\rho\in\mathbb R$, let $S$ be a subset of
$\{1,\ldots,n\}$, and let $g$ be any real-valued function on its subsets. The weighted
average of $g(S\cap K)$ over all subsets $K\subseteq\{1,\ldots,n\}$, with weight
$\rho^{|K|}(1-\rho)^{n-|K|}$, equals $\sum_{J\subseteq S}\rho^{|J|}(1-\rho)^{|S\setminus
J|}g(J)$.
\end{theorem}

\begin{definition}[Signed random-subset Fourier average]
\lean{FABL.oxrRandomSubsetFourierAverage}
\statementsource{Analysis-of-Boolean-Functions-by-Ryan-ODonnell}{pdf-ffbe9c7b3306-p201-b003}
\label{FABL.oxrRandomSubsetFourierAverage}
\leanok
\uses{FABL.expectDeltaRandomSubset, FABL.fourierCoeff}
For a natural number $n$, a real parameter $\delta$, a real-valued function $f$ on the
$n$-dimensional sign cube, and a subset $S$ of the coordinate indices, define the
average over random subsets $K$ with parameter $1-\delta$ of $(-1)^{|S\cap K|}\widehat
f(S\cap K)$, where $\widehat f$ denotes the Fourier coefficient of $f$.
\end{definition}

\begin{definition}[Absolute random-subset Fourier average]
\lean{FABL.oxrRandomSubsetAbsFourierAverage}
\statementsource{Analysis-of-Boolean-Functions-by-Ryan-ODonnell}{pdf-ffbe9c7b3306-p201-b003}
\label{FABL.oxrRandomSubsetAbsFourierAverage}
\leanok
\uses{FABL.expectDeltaRandomSubset, FABL.fourierCoeff}
For a natural number $n$, a real parameter $\delta$, a real-valued function $f$ on the
$n$-dimensional sign cube, and a subset $S$ of the coordinate indices, define the
average over random subsets $K$ with parameter $1-\delta$ of $|\widehat f(S\cap K)|$,
where $\widehat f$ denotes the Fourier coefficient of $f$.
\end{definition}

\begin{theorem}[Absolute expectation bounded by the second moment]
\lean{FABL.expectDeltaRandomSubset_abs_le_sqrt_sq}
\label{FABL.expectDeltaRandomSubset_abs_le_sqrt_sq}
\leanok
\uses{FABL.deltaRandomSubsetWeight, FABL.deltaRandomSubsetWeight_nonneg, FABL.expectDeltaRandomSubset, FABL.sum_deltaRandomSubsetWeight}
\difficulty{4}
Let $n$ be a nonnegative integer, let $\rho\in[0,1]$, and let $g$ be any real-valued
function on the subsets of $\{1,\ldots,n\}$. The weighted average of $|g(J)|$ over
subsets $J$, with each subset given weight $\rho^{|J|}(1-\rho)^{n-|J|}$, is at most the
square root of the weighted average of $g(J)^2$ with the same weights.
\end{theorem}

\begin{theorem}[Fourier average bound above a set-size threshold]
\lean{FABL.oxrRandomSubsetAbsFourierAverage_le_sqrt_of_lt_card}
\label{FABL.oxrRandomSubsetAbsFourierAverage_le_sqrt_of_lt_card}
\leanok
\uses{FABL.BoundedCubeFunction, FABL.BoundedCubeFunction.sum_sq_fourierCoeff_le_one, FABL.expectDeltaRandomSubset, FABL.expectDeltaRandomSubset_abs_le_sqrt_sq, FABL.expectDeltaRandomSubset_intersection_eq_sum, FABL.fourierCoeff, FABL.oxrRandomSubsetAbsFourierAverage}
\difficulty{5}
Let $n$ be a natural number, let $\delta\in[0,1]$ with $\delta\le\tfrac12$, and let $f$
be a real-valued function on the $n$-dimensional Boolean cube whose values lie in
$[-1,1]$. If $t$ is a natural number and $S$ is a set of coordinates with $t<|S|$, then
the expected absolute Fourier coefficient indexed by $S\cap K$, where $K$ is sampled
with the subset weights at parameter $1-\delta$, is at most $\sqrt{(1-\delta)^t}$.
\end{theorem}

\begin{theorem}[A signed Fourier average bound]
\lean{FABL.neg_oxrRandomSubsetFourierAverage_le_abs}
\label{FABL.neg_oxrRandomSubsetFourierAverage_le_abs}
\leanok
\uses{FABL.deltaRandomSubsetWeight, FABL.deltaRandomSubsetWeight_nonneg, FABL.expectDeltaRandomSubset, FABL.fourierCoeff, FABL.oxrRandomSubsetAbsFourierAverage, FABL.oxrRandomSubsetFourierAverage}
\difficulty{3}
Let $n$ be a natural number, let $0\le\delta\le1$, let $f$ be any real-valued function
on the $n$-dimensional sign cube, and let $S$ be a subset of its coordinates. For a
random subset $K$ formed by including each coordinate independently with probability
$1-\delta$, the negative of the average of $(-1)^{|S\cap K|}\widehat f(S\cap K)$ is at
most the average of $|\widehat f(S\cap K)|$, where $\widehat f$ denotes the uniform
Fourier coefficient.
\end{theorem}

\begin{theorem}[Vanishing even Fourier coefficients]
\lean{FABL.fourierCoeff_even_eq_zero_of_odd}
\proofsource{Analysis-of-Boolean-Functions-by-Ryan-ODonnell}{
  pdf-ffbe9c7b3306-p035-b008,
  pdf-ffbe9c7b3306-p035-b009,
  pdf-ffbe9c7b3306-p122-b004
}
\label{FABL.fourierCoeff_even_eq_zero_of_odd}
\leanok
\uses{FABL.fourierCoeff, FABL.fourierCoeff_oddPart_of_even, FABL.oddPart}
\difficulty{3}
Let $f$ be a real-valued function on the $n$-dimensional sign cube satisfying
$f(-x)=-f(x)$ for every $x$, and let $S$ be a set of coordinates of even cardinality.
Then the Fourier coefficient of $f$ at $S$ is zero.
\end{theorem}

\begin{theorem}[Neg noise stability neg le of odd]
\lean{FABL.neg_noiseStability_neg_le_of_odd}
\statementsource{Analysis-of-Boolean-Functions-by-Ryan-ODonnell}{pdf-ffbe9c7b3306-p064-b002}
\label{FABL.neg_noiseStability_neg_le_of_odd}
\leanok
\uses{FABL.BoundedCubeFunction, FABL.BoundedCubeFunction.sum_sq_fourierCoeff_le_one, FABL.fourierCoeff, FABL.fourierCoeff_even_eq_zero_of_odd, FABL.noiseStability, FABL.noiseStability_eq_sum_rho_pow_mul_sq_fourierCoeff}
\difficulty{5}
% dataset-pipeline: stub -- prose pending build_informal_statements.py
Statement pending regeneration.
\end{theorem}

\begin{theorem}[A bound for a signed Fourier sum]
\lean{FABL.neg_oxrFourierSum_le_low_add_tail}
\label{FABL.neg_oxrFourierSum_le_low_add_tail}
\leanok
\uses{FABL.BoundedCubeFunction, FABL.BoundedCubeFunction.sum_sq_fourierCoeff_le_one, FABL.fourierCoeff, FABL.neg_oxrRandomSubsetFourierAverage_le_abs, FABL.oxrRandomSubsetAbsFourierAverage, FABL.oxrRandomSubsetAbsFourierAverage_le_sqrt_of_lt_card, FABL.oxrRandomSubsetFourierAverage}
\difficulty{5}
Let $n,t$ be natural numbers, let $0\le\delta\le1$ with $\delta\le1/2$, and let $f$ be a
real-valued function on the $n$-dimensional sign cube whose values lie in $[-1,1]$. For
each coordinate subset $S\subseteq\{1,\ldots,n\}$, write $\widehat f(S)$ for its uniform
Fourier coefficient, and let $A_\delta(S)$ and $B_\delta(S)$ be the averages over random
subsets $K$, in which each coordinate is included independently with probability
$1-\delta$, of $(-1)^{|S\cap K|}\widehat f(S\cap K)$ and $|\widehat f(S\cap K)|$,
respectively. Then $-\sum_S \widehat f(S)^2 A_\delta(S)\le
\sqrt{(1-\delta)^t}+\sum_{S:\,|S|\le t}\widehat f(S)^2 B_\delta(S)$.
\end{theorem}

\begin{theorem}[Odd-function Fourier term as a subset average]
\lean{FABL.oxrTripleFourierTerm_eq_randomSubsetAverage_of_odd}
\label{FABL.oxrTripleFourierTerm_eq_randomSubsetAverage_of_odd}
\leanok
\uses{FABL.expectDeltaRandomSubset, FABL.expectDeltaRandomSubset_intersection_eq_sum, FABL.fourierCoeff, FABL.fourierCoeff_even_eq_zero_of_odd, FABL.oxrRandomSubsetFourierAverage}
\difficulty{5}
Let $n$ be a natural number, let $\delta$ be real, and let $f$ be an odd real-valued
function on the $n$-dimensional Boolean cube. For every set of coordinates $S$,
\[
\widehat f(S)^2\sum_{J\subseteq S} (-(1-\delta))^{|J|}(-\delta)^{|S|-|J|}\widehat f(J)
=\widehat f(S)^2\,\mathbb E_K\!\left[(-1)^{|S\cap K|}\widehat f(S\cap K)\right],
\]
where $\widehat f$ denotes the Fourier coefficient and $K$ is sampled with the subset
weights at parameter $1-\delta$.
\end{theorem}

\begin{theorem}[Triple-product expectation as a subset average]
\lean{FABL.pmfExpectation_oxrTripleCubePMF_tripleProduct_eq_randomSubsetAverage}
\label{FABL.pmfExpectation_oxrTripleCubePMF_tripleProduct_eq_randomSubsetAverage}
\leanok
\uses{FABL.fourierCoeff, FABL.oxrQueryPoint, FABL.oxrRandomSubsetFourierAverage, FABL.oxrTripleCubePMF, FABL.oxrTripleFourierTerm_eq_randomSubsetAverage_of_odd, FABL.pmfExpectation, FABL.pmfExpectation_oxrTripleCubePMF_tripleProduct_fourier}
\difficulty{3}
Let $n$ be a natural number, let $\delta\in[0,1]$, and let $f:\{-1,1\}^n\to\mathbb R$ be
odd, so that $f(-x)=-f(x)$. Choose independently at each coordinate a sign triple by
selecting $(U,V,-V)$ with probability $\delta$ and $(U,V,-UV)$ otherwise, with
independent uniform signs $U,V$. The expectation of the product of $f$ on the three
query vectors is $\sum_{S\subseteq\{1,\ldots,n\}}\widehat f(S)^2 A_{\delta,f}(S)$, where
$\widehat f(S)$ is the Fourier coefficient of $f$ at $S$ and $A_{\delta,f}(S)$ is the
average over subsets $K$ with weight $(1-\delta)^{|K|}\delta^{n-|K|}$ of $(-1)^{|S\cap
K|}\widehat f(S\cap K)$.
\end{theorem}

\begin{definition}[OXR acceptance probability]
\lean{FABL.oxrAcceptanceProbability}
\statementsource{Analysis-of-Boolean-Functions-by-Ryan-ODonnell}{pdf-ffbe9c7b3306-p201-b003}
\label{FABL.oxrAcceptanceProbability}
\leanok
\uses{FABL.BoundedCubeFunction, FABL.SignPredicate.randomizedValue, FABL.oxrPredicate, FABL.oxrQueryPoint, FABL.oxrTripleCubePMF, FABL.pmfExpectation}
For a natural number $n$, let $\delta\in[0,1]$ and let $f$ be a function from the
$n$-dimensional sign cube to $[-1,1]$. Define its acceptance probability to be the
expectation, over the correlated triple-cube sample with parameter $\delta$, of the OXR
predicate's randomized value on the three function values at the sample's query points.
\end{definition}

\begin{theorem}[Odd tester acceptance identity]
\lean{FABL.oxrOddSignTester_acceptanceProbability}
\label{FABL.oxrOddSignTester_acceptanceProbability}
\leanok
\uses{FABL.BoundedCubeFunction, FABL.BoundedCubeFunction.oddPart, FABL.OXRSignTester.toSignTester, FABL.Sign, FABL.SignQueryCheck, FABL.SignQueryCheck.oddSymmetrizedValue, FABL.SignQueryCheck.oddSymmetrizedValue_eq_oddPart, FABL.SignTester.acceptanceProbability, FABL.oxrAcceptanceProbability, FABL.oxrOddSignTester, FABL.oxrPredicate, FABL.oxrQueryPoint, FABL.oxrTripleCubePMF, FABL.pmfExpectation, FABL.pmfExpectation_bind, FABL.pmfExpectation_map, FABL.pmfExpectation_uniformPMF_eq_expect}
\difficulty{3}
Let $n\ge0$, let $\delta\in[0,1]$, and let $f:\{\pm1\}^n\to[-1,1]$. The acceptance
probability of the odd OXR tester with parameter $\delta$ equals the acceptance
probability of the three-query OXR test with parameter $\delta$ applied to the odd part
$x\mapsto(f(x)-f(-x))/2$.
\end{theorem}

\begin{theorem}[OXR acceptance identity]
\lean{FABL.oxrAcceptanceProbability_eq_equationSevenSeven}
\label{FABL.oxrAcceptanceProbability_eq_equationSevenSeven}
\leanok
\uses{FABL.BoundedCubeFunction, FABL.Sign, FABL.fourierCoeff, FABL.fourierCoeff_even_eq_zero_of_odd, FABL.noiseStability, FABL.oxrAcceptanceProbability, FABL.oxrPredicate_randomizedValue, FABL.oxrQueryPoint, FABL.oxrRandomSubsetFourierAverage, FABL.pmfExpectation_add, FABL.pmfExpectation_const, FABL.pmfExpectation_const_mul, FABL.pmfExpectation_oxrTripleCubePMF_first, FABL.pmfExpectation_oxrTripleCubePMF_second_mul_third, FABL.pmfExpectation_oxrTripleCubePMF_tripleProduct_eq_randomSubsetAverage}
\difficulty{4}
Let $n\ge 0$, let $\delta\in[0,1]$, and let $f:\{\pm1\}^n\to[-1,1]$ be odd. For each
$S\subseteq[n]$, write $\widehat f(S)$ for its Fourier coefficient and set
$A_\delta(S)=\sum_{K\subseteq[n]}(1-\delta)^{|K|}\delta^{n-|K|}(-1)^{|S\cap K|}\widehat
f(S\cap K)$. Then the acceptance probability of the three-query OXR test with parameter
$\delta$ is
$\frac34-\frac14\operatorname{Stab}_{-\delta}(f)-\frac14\sum_{S\subseteq[n]}\widehat
f(S)^2 A_\delta(S)$, where $\operatorname{Stab}_{-\delta}(f)$ is the noise stability at
correlation $-\delta$.
\end{theorem}

\begin{theorem}[Oxr acceptance probability le exercise seven twenty nine]
\lean{FABL.oxrAcceptanceProbability_le_exerciseSevenTwentyNine}
\statementsource{Analysis-of-Boolean-Functions-by-Ryan-ODonnell}{
  pdf-ffbe9c7b3306-p201-b003,
  pdf-ffbe9c7b3306-p201-b006
}
\label{FABL.oxrAcceptanceProbability_le_exerciseSevenTwentyNine}
\leanok
\uses{FABL.BoundedCubeFunction, FABL.fourierCoeff, FABL.neg_noiseStability_neg_le_of_odd, FABL.neg_oxrFourierSum_le_low_add_tail, FABL.oxrAcceptanceProbability, FABL.oxrAcceptanceProbability_eq_equationSevenSeven, FABL.oxrRandomSubsetAbsFourierAverage}
\difficulty{3}
% dataset-pipeline: stub -- prose pending build_informal_statements.py
Statement pending regeneration.
\end{theorem}

\begin{theorem}[Abs fourier coeff le sqrt of small stable influences]
\lean{FABL.abs_fourierCoeff_le_sqrt_of_small_stableInfluences}
\statementsource{Analysis-of-Boolean-Functions-by-Ryan-ODonnell}{pdf-ffbe9c7b3306-p195-b001}
\label{FABL.abs_fourierCoeff_le_sqrt_of_small_stableInfluences}
\leanok
\uses{FABL.HasSmallStableInfluences, FABL.fourierCoeff, FABL.stableInfluence}
\difficulty{5}
% dataset-pipeline: stub -- prose pending build_informal_statements.py
Statement pending regeneration.
\end{theorem}

\begin{theorem}[Fourier average under low-degree regularity]
\lean{FABL.oxrRandomSubsetAbsFourierAverage_le_of_lowDegreeRegular}
\label{FABL.oxrRandomSubsetAbsFourierAverage_le_of_lowDegreeRegular}
\leanok
\uses{FABL.IsLowDegreeFourierRegular, FABL.deltaRandomSubsetWeight, FABL.deltaRandomSubsetWeight_nonneg, FABL.expectDeltaRandomSubset, FABL.fourierCoeff, FABL.fourierCoeff_even_eq_zero_of_odd, FABL.oxrRandomSubsetAbsFourierAverage, FABL.sum_deltaRandomSubsetWeight}
\difficulty{4}
Let $n$ be a natural number, let $\delta\in[0,1]$, and let $f$ be an odd real-valued
function on the $n$-dimensional Boolean cube. Suppose $R\ge 0$ and every nonempty
Fourier coefficient of $f$ indexed by a set of at most $t$ coordinates has absolute
value at most $R$. Then, for every set $S$ of at most $t$ coordinates, the expected
absolute Fourier coefficient indexed by $S\cap K$, where $K$ is sampled with the subset
weights at parameter $1-\delta$, is at most $R$.
\end{theorem}

\begin{theorem}[Low-degree Fourier remainder bound]
\lean{FABL.oxrLowDegreeRemainder_le_of_lowDegreeRegular}
\label{FABL.oxrLowDegreeRemainder_le_of_lowDegreeRegular}
\leanok
\uses{FABL.BoundedCubeFunction, FABL.BoundedCubeFunction.sum_sq_fourierCoeff_le_one, FABL.IsLowDegreeFourierRegular, FABL.fourierCoeff, FABL.oxrRandomSubsetAbsFourierAverage, FABL.oxrRandomSubsetAbsFourierAverage_le_of_lowDegreeRegular}
\difficulty{4}
Let $n\ge0$, let $\delta\in[0,1]$, and let $f:\{\pm1\}^n\to[-1,1]$ be odd. Suppose
$R\ge0$, $t\in\mathbb N$, and every nonempty $S\subseteq[n]$ with $|S|\le t$ has Fourier
coefficient satisfying $|\widehat f(S)|\le R$. For $S\subseteq[n]$, let $B_\delta(S)$ be
the average of $|\widehat f(S\cap K)|$ over subsets $K\subseteq[n]$ with weights
$(1-\delta)^{|K|}\delta^{n-|K|}$. Then $\sum_{S\subseteq[n],\,|S|\le t}\widehat f(S)^2
B_\delta(S)\le R$.
\end{theorem}

\begin{theorem}[Dictator acceptance probability]
\lean{FABL.oxrAcceptanceProbability_dictator}
\label{FABL.oxrAcceptanceProbability_dictator}
\leanok
\uses{FABL.BooleanFunction.toReal, FABL.BoundedCubeFunction.ofBoolean, FABL.dictator, FABL.dictator_toReal_eq_monomial_singleton, FABL.expectDeltaRandomSubset, FABL.expectDeltaRandomSubset_intersection_eq_sum, FABL.fourierCoeff, FABL.fourierCoeff_monomial, FABL.monomial, FABL.noiseStability, FABL.noiseStability_dictator, FABL.oxrAcceptanceProbability, FABL.oxrAcceptanceProbability_eq_equationSevenSeven, FABL.oxrRandomSubsetFourierAverage, FABL.signValue, FABL.signValue_neg}
\difficulty{5}
Let $n$ be a natural number, let $\delta\in[0,1]$, and choose a coordinate $i$ of the
$n$-dimensional sign cube. The acceptance probability of the specified OXR test on the
dictator function $x\mapsto x_i$ is $1$.
\end{theorem}

\begin{definition}[Håstad three-SAT degree parameter]
\lean{FABL.hastadThreeSatDegree}
\label{FABL.hastadThreeSatDegree}
\leanok
For $\delta\in(0,1/8)$, define the degree parameter to be a natural number $k$ chosen
so that $(1-4\delta)^k<(4\delta)^2$.
\end{definition}

\begin{theorem}[Three-SAT degree tail bound]
\lean{FABL.hastadThreeSatDegree_tail_lt}
\label{FABL.hastadThreeSatDegree_tail_lt}
\leanok
\uses{FABL.hastadThreeSatDegree}
\difficulty{3}
For every $\delta\in(0,1/8)$, let $d$ be the degree selected by the definition of the
three-SAT degree. Then $\sqrt{(1-4\delta)^d}<4\delta$.
\end{theorem}

\begin{definition}[Håstad three-SAT OXR error bound]
\lean{FABL.hastadThreeSatOXRError}
\label{FABL.hastadThreeSatOXRError}
\leanok
For a natural number $t$ and a real number $\varepsilon$, define the OXR error to be
$\min\{1,\frac14\sqrt{\varepsilon/(1-\varepsilon)^{t-1}}\}$ when $\varepsilon<1$, and to
be $1$ when $\varepsilon\ge1$.
\end{definition}

\begin{theorem}[Three-SAT OXR error range]
\lean{FABL.hastadThreeSatOXRError_mem_Icc}
\label{FABL.hastadThreeSatOXRError_mem_Icc}
\leanok
\uses{FABL.hastadThreeSatOXRError}
\difficulty{4}
For every natural number $t$ and every $\varepsilon\in[0,1]$, the three-SAT OXR error at
parameters $t$ and $\varepsilon$ lies in $[0,1]$.
\end{theorem}

\begin{theorem}[Vanishing three-SAT OXR error]
\lean{FABL.hastadThreeSatOXRError_tendsto_zero}
\label{FABL.hastadThreeSatOXRError_tendsto_zero}
\leanok
\uses{FABL.hastadThreeSatOXRError}
\difficulty{4}
For each natural number $t$, the function $\\varepsilon \\mapsto
\\mathrm{hastadThreeSatOXRError}(t,\\varepsilon)$ tends to $0$ as $\\varepsilon$ tends
to $0$ through positive real values.
\end{theorem}

\begin{definition}[Håstad OXR test family]
\lean{FABL.hastadOXRTestFamily}
\statementsource{Analysis-of-Boolean-Functions-by-Ryan-ODonnell}{pdf-ffbe9c7b3306-p201-b003}
\label{FABL.hastadOXRTestFamily}
\leanok
\uses{FABL.OXRSignTester, FABL.oxrOddSignTester}
For a real number $\delta$ with $0<\delta<1/8$ and a positive natural number $n$, let
the Håstad OXR test family be the OXR sign tester on $n$ coordinates obtained by
applying the odd sign tester with parameter $4\delta$.
\end{definition}

\begin{theorem}[OXR dictator versus no-notables test]
\lean{FABL.hastadOXR_dictatorVsNoNotables}
\proofsource{Analysis-of-Boolean-Functions-by-Ryan-ODonnell}{
  pdf-ffbe9c7b3306-p201-b003,
  pdf-ffbe9c7b3306-p201-b006
}
\label{FABL.hastadOXR_dictatorVsNoNotables}
\leanok
\uses{FABL.BoundedCubeFunction.oddPart, FABL.BoundedCubeFunction.oddPart_ofBoolean_dictator, FABL.BoundedCubeFunction.ofBoolean, FABL.HasSmallStableInfluences, FABL.IsDictatorVsNoNotablesTest, FABL.IsLowDegreeFourierRegular, FABL.OXRSignTester.toSignTester, FABL.OXRSignTester.toSignTester_usesMaxOXR, FABL.SignTester.acceptanceProbability, FABL.SignTester.acceptanceProbability_mem_Icc, FABL.abs_fourierCoeff_le_sqrt_of_small_stableInfluences, FABL.dictator, FABL.hastadOXRTestFamily, FABL.hastadThreeSatDegree, FABL.hastadThreeSatDegree_tail_lt, FABL.hastadThreeSatOXRError, FABL.hastadThreeSatOXRError_mem_Icc, FABL.hastadThreeSatOXRError_tendsto_zero, FABL.maxOXRPredicates, FABL.maxOXRPredicates_finite, FABL.odd_oddPart, FABL.oxrAcceptanceProbability_dictator, FABL.oxrAcceptanceProbability_le_exerciseSevenTwentyNine, FABL.oxrLowDegreeRemainder_le_of_lowDegreeRegular, FABL.oxrOddSignTester, FABL.oxrOddSignTester_acceptanceProbability, FABL.stableInfluence_oddPart_le}
\difficulty{5}
For each real number $\delta$ with $0<\delta<1/8$, the test family obtained from the odd
OXR sign tester with parameter $4\delta$ is a dictator-versus-no-notables test for the
family of all input-negations of the three-input OXR predicate. Its parameters are
$\alpha=3/4+2\delta$ and $\beta=1$, and its error function is the one determined by the
degree chosen from $\delta$.
\end{theorem}

\begin{definition}[Håstad three-SAT tester family]
\lean{FABL.hastadThreeSatTesterFamily}
\statementsource{Analysis-of-Boolean-Functions-by-Ryan-ODonnell}{pdf-ffbe9c7b3306-p201-b003}
\label{FABL.hastadThreeSatTesterFamily}
\leanok
\uses{FABL.OXRSignTester.toE3SatTester, FABL.SignTester, FABL.hastadOXRTestFamily}
For each real number $\delta$ with $0<\delta<1/8$ and each positive integer $n$, define
a sign tester with input length $n$ by converting the Håstad OXR test family with
parameters $\delta$ and $n$ into a three-SAT tester.
\end{definition}

\begin{theorem}[Three-SAT dictator versus no notable coordinates]
\lean{FABL.hastadThreeSat_dictatorVsNoNotables}
\proofsource{Analysis-of-Boolean-Functions-by-Ryan-ODonnell}{pdf-ffbe9c7b3306-p201-b003}
\label{FABL.hastadThreeSat_dictatorVsNoNotables}
\leanok
\uses{FABL.IsDictatorVsNoNotablesTest, FABL.hastadOXRTestFamily, FABL.hastadOXR_dictatorVsNoNotables, FABL.hastadThreeSatDegree, FABL.hastadThreeSatOXRError, FABL.hastadThreeSatTesterFamily, FABL.maxE3SatSignPredicates, FABL.oxrDictatorTest_to_e3Sat_affine, FABL.oxrE3SatReplacementError}
\difficulty{2}
For every $\delta\in(0,1/8)$, there is a family of finite-seed sign testers for
three-literal sign clauses with the following properties. The predicate set is finite;
the thresholds are $\alpha=7/8+\delta$ and $\beta=1$; and the error function is half the
three-SAT OXR error evaluated at the degree selected for $\delta$. This error takes
values in $[0,1]$ and tends to zero as its argument approaches zero through positive
values. Every test uses an allowed predicate, every dictator is accepted with
probability at least $\beta$, and every function whose stable influences at parameters
$(\varepsilon,\varepsilon)$ are all at most $\varepsilon$ is accepted with probability
at most $\alpha$ plus the error at $\varepsilon$, for every dimension, every positive
dimension, and every $\varepsilon\in[0,1]$.
\end{theorem}
